% !TEX program = xelatex
\documentclass[12pt,a4paper]{ctexart}

% ══════════════════════════════════════════════════════════════════════
% PACKAGES
% ══════════════════════════════════════════════════════════════════════
\usepackage{amsmath,amssymb,amsthm}
\usepackage{mathtools}
\usepackage{bm}
\usepackage{graphicx}
\usepackage{hyperref}
\usepackage{geometry}
\usepackage{enumitem}
\usepackage{booktabs}
\usepackage{tikz-cd}
\usepackage{xcolor}
\usepackage{cleveref}
\usepackage{bbm}
\usepackage{physics}
\usepackage{array}
\usepackage{multirow}
\usepackage{makecell}
\usepackage{longtable}

\geometry{margin=2.5cm}

\hypersetup{
    colorlinks=true,
    linkcolor=blue,
    citecolor=blue,
    urlcolor=blue,
    pdftitle={SCX Gauge Theory Formalized: O(d) Lattice Gauge, Dijkgraaf-Witten TQFT, and Information-Geometric Bulk-Boundary},
    pdfauthor={SCX},
    pdfsubject={SCX Gauge Theory},
}

% ══════════════════════════════════════════════════════════════════════
% THEOREM ENVIRONMENTS
% ══════════════════════════════════════════════════════════════════════
\theoremstyle{definition}
\newtheorem{definition}{定义}[section]
\newtheorem{theorem}{定理}[section]
\newtheorem{proposition}{命题}[section]
\newtheorem{corollary}{推论}[section]
\newtheorem{lemma}{引理}[section]
\newtheorem{conjecture}{猜想}[section]
\newtheorem{problem}{问题}[section]

\theoremstyle{remark}
\newtheorem{remark}{注记}[section]
\newtheorem{example}{例}[section]

\theoremstyle{definition}
\newtheorem*{definition*}{Definition}
\newtheorem*{theorem*}{Theorem}
\newtheorem*{proposition*}{Proposition}
\newtheorem*{corollary*}{Corollary}
\newtheorem*{lemma*}{Lemma}
\newtheorem*{remark*}{Remark}
\newtheorem*{problem*}{Open Problem}

% ══════════════════════════════════════════════════════════════════════
% CUSTOM COMMANDS — ALL \newcommand DEFINED HERE
% ══════════════════════════════════════════════════════════════════════

% ── Standard math blackboard ──
\newcommand{\R}{\mathbb{R}}
\newcommand{\C}{\mathbb{C}}
\newcommand{\Z}{\mathbb{Z}}
\newcommand{\N}{\mathbb{N}}
\newcommand{\Q}{\mathbb{Q}}
\newcommand{\F}{\mathbb{F}}
\newcommand{\E}{\mathbb{E}}
\newcommand{\I}{\mathbb{I}}

% ── Groups ──
\newcommand{\Od}{O(d)}
\newcommand{\SOd}{SO(d)}
\newcommand{\Ud}{U(d)}
\newcommand{\SUd}{SU(d)}
\newcommand{\SEN}{SE(d)}
\newcommand{\GLd}{GL(d,\R)}

% ── Graph notation ──
\newcommand{\grph}{\mathcal{G}}
\newcommand{\verts}{\mathcal{V}}
\newcommand{\edgs}{\mathcal{E}}
\newcommand{\faces}{\mathcal{F}}

\theoremstyle{definition}
\newtheorem*{definition*}{Definition}
\newtheorem*{theorem*}{Theorem}
\newtheorem*{proposition*}{Proposition}
\newtheorem*{corollary*}{Corollary}
\newtheorem*{lemma*}{Lemma}
\newtheorem*{remark*}{Remark}
\newtheorem*{problem*}{Open Problem}
% ══════════════════════════════════════════════════════════════════════
% ALL \newcommand DEFINED HERE
% ══════════════════════════════════════════════════════════════════════

% ── Standard math blackboard ──
\newcommand{\R}{\mathbb{R}}
\newcommand{\C}{\mathbb{C}}
\newcommand{\Z}{\mathbb{Z}}
\newcommand{\N}{\mathbb{N}}
\newcommand{\Q}{\mathbb{Q}}
\newcommand{\F}{\mathbb{F}}
\newcommand{\E}{\mathbb{E}}
\newcommand{\I}{\mathbb{I}}

% ── Groups ──
\newcommand{\Od}{O(d)}
\newcommand{\SOd}{SO(d)}
\newcommand{\Ud}{U(d)}
\newcommand{\SUd}{SU(d)}
\newcommand{\SEN}{SE(d)}
\newcommand{\GLd}{GL(d,\R)}
\newcommand{\SOthree}{SO(3)}

% ── Graph notation ──
\newcommand{\grph}{\mathcal{G}}
\newcommand{\verts}{\mathcal{V}}
\newcommand{\edgs}{\mathcal{E}}
\newcommand{\faces}{\mathcal{F}}
\newcommand{\Loops}{\mathcal{L}}

% ── Discrete Hodge operators ──
\newcommand{\dzero}{d_0}
\newcommand{\done}{d_1}
\newcommand{\dtwo}{d_2}
\newcommand{\dzeroT}{d_0^\dagger}
\newcommand{\doneT}{d_1^\dagger}
\newcommand{\dtwoT}{d_2^\dagger}
\newcommand{\Lzero}{\Delta_0}
\newcommand{\Lone}{\Delta_1}
\newcommand{\Ltwo}{\Delta_2}

% ── Gauge theory ──
\newcommand{\curv}{\Omega}
\newcommand{\Hol}{\mathrm{Hol}}
\newcommand{\triv}{\mathrm{triv}}
\newcommand{\Ad}{\mathrm{Ad}}
\newcommand{\cercis}{\mathcal{C}}
\newcommand{\cercisnorm}{\bar{\mathcal{C}}}
\newcommand{\SCX}{\textsc{SCX}}
\newcommand{\Situs}{\textsc{Situs}}
\newcommand{\Spring}{\textsc{Spring}}
\newcommand{\Yajie}{\textsc{Yajie}}
\newcommand{\ACE}{\textsc{ACE}}
\newcommand{\MILP}{\textsc{MILP-gf}}

% ── Information geometry ──
\newcommand{\Fisher}{\mathcal{I}}
\newcommand{\KL}{\mathrm{KL}}
\newcommand{\Dist}{\mathcal{D}}
\newcommand{\ExpFam}{\mathcal{E}}
\newcommand{\GeoDist}{d_{\mathcal{M}}}

% ── TQFT ──
\newcommand{\DW}{\mathrm{DW}}
\newcommand{\ZBetti}{\beta_1}
\newcommand{\ZBettiTwo}{\beta_2}
\newcommand{\Gflat}{\mathcal{M}_{\mathrm{flat}}}
\newcommand{\Gflatquot}{\mathcal{M}_{\mathrm{flat}}/\mathcal{G}}

% ── Operators ──
\DeclareMathOperator{\diag}{diag}
\DeclareMathOperator{\sgn}{sgn}
\DeclareMathOperator{\im}{im}
\DeclareMathOperator{\coker}{coker}
\DeclareMathOperator{\Hom}{Hom}
\DeclareMathOperator{\Aut}{Aut}
\DeclareMathOperator{\vol}{vol}
\DeclareMathOperator{\Spec}{Spec}
\DeclareMathOperator{\rk}{rank}
\DeclareMathOperator{\tr}{tr}
\DeclareMathOperator{\Rep}{Rep}
\DeclareMathOperator{\Irr}{Irr}
\DeclareMathOperator{\Span}{span}
\DeclareMathOperator{\dist}{dist}
\DeclareMathOperator{\grad}{grad}
\DeclareMathOperator{\Log}{Log}

% ── DW-specific ──
\newcommand{\ZofG}{Z(G)}
\newcommand{\ClG}{\mathrm{Cl}(G)}
\newcommand{\Irrep}{\widehat{G}}

% ── Bundle notation ──
\newcommand{\bun}{P}
\newcommand{\base}{\mathcal{X}}

% ══════════════════════════════════════════════════════════════════════
% TITLE
% ══════════════════════════════════════════════════════════════════════
\title{
    {\Huge \textbf{SCX规范理论的形式化完善}}\\[0.5em]
    {\Large Formal Completion of SCX Gauge Theory}\\[1.0em]
    {\normalsize $\Od$格点规范·Dijkgraaf-Witten离散TQFT·信息几何体-边界对应}\\[0.3em]
    {\normalsize $\Od$ Lattice Gauge · Dijkgraaf-Witten Discrete TQFT · Information-Geometric Bulk-Boundary}
}

\author{SCX 理论架构组 \\ SCX Theory Architecture Group}
\date{\today}

\begin{document}
\maketitle

\begin{center}
\footnotesize
\textbf{版本：} v2.0 — 形式化完善版 \quad | \quad
\textbf{状态：} 预印本 \quad | \quad
\textbf{分类：} SCX理论体系 — 规范场论卷·形式化篇
\end{center}

% ══════════════════════════════════════════════════════════════════════
% ABSTRACT (Bilingual)
% ══════════════════════════════════════════════════════════════════════
\begin{abstract}

\noindent\textbf{中文摘要：}
本文完成SCX多专家系统规范理论在三个相互关联的领域中的严格数学形式化，超越此前图霍奇理论的阿贝尔平移规范框架。
(1)~\textbf{$\Od$格点规范理论：}将专家的旋转对称性形式化为图上的非阿贝尔$\Od$格点规范理论——顶点承载$\Od$参考标架，
边承载$\Od$值ACE联络势，规范变换为非阿贝尔共轭$A_e \mapsto g_v \circ A_e \circ g_u^{-1}$，
曲率定义为面上Wilson圈$\Hol(\face) = \prod_{e \in \partial \face} A_e^{\sigma_e}$。
我们证明在小联络极限下，$\Od$规范固定等价于李代数$\so(d)$上的标准图霍奇分解，
且在大扭曲情形下解的存在性由$\Od^{|\verts|}$的紧致性保证。
(2)~\textbf{Dijkgraaf-Witten离散TQFT：}将专家图构造为2-复形（顶点-边-面），
在其上定义Dijkgraaf-Witten离散拓扑量子场论。对于有限群$G$，配分函数
$Z_{\DW}(\mathcal{K}, G) = |\Hom(\pi_1(\mathcal{K}), G)/G|$精确计数规范不等价的平坦联络。
我们证明SCX 2-复形同伦于$M$顶点完全图$K_M$，其第一Betti数为$\ZBetti = \binom{M}{2} - M + 1 = (M-1)(M-2)/2$，
且Cercis分数的调和分量$r_{\text{harm}} \in \ker(\Lone)$正是平坦联络模空间$\Gflatquot$在平凡联络处的切空间。
(3)~\textbf{信息几何体-边界对应：}在$\Situs$流形上赋予Fisher信息度量$\Fisher_{ij}(\theta)$，
将专家输出分布建模为指数族$\ExpFam$。我们证明在此族下，Fisher测地线长度与KL散度满足局部等价性：
$\GeoDist(p,q)^2 = 2\KL(p\|q) + O(\|p-q\|^3)$。
进一步证明Cercis分数在指数族近似下等于边界观测分布到调和体态的最小Fisher测地距离：
$\cercis = \min_{h \in \ker(\Lone)} \GeoDist(P_{\text{obs}}, P_h)^2$。
本文所有概念均以定理、引理、推论或定义的形式呈现，无类比语言。
第\ref{sec:comparison}节详细比较本文形式化与旧有fiber\_bundle.tex之间的差异。

\vspace{0.8em}
\noindent\textbf{English Abstract:}
This paper completes the rigorous mathematical formalization of SCX multi-expert gauge theory in three interconnected domains,
transcending the abelian translational gauge framework of prior graph Hodge theory.
(1)~\textbf{$\Od$ Lattice Gauge Theory:} Expert rotational symmetry is formalized as non-abelian $\Od$ lattice gauge theory on a graph —
vertices carry $\Od$ reference frames, edges carry $\Od$-valued ACE connection potentials,
gauge transformations are non-abelian conjugation $A_e \mapsto g_v \circ A_e \circ g_u^{-1}$,
and curvature is the face Wilson loop $\Hol(\face) = \prod_{e \in \partial \face} A_e^{\sigma_e}$.
We prove that in the small-connection limit, $\Od$ gauge-fixing reduces to standard graph Hodge decomposition on the Lie algebra $\so(d)$,
and in the large-distortion regime, existence is guaranteed by compactness of $\Od^{|\verts|}$.
(2)~\textbf{Dijkgraaf-Witten Discrete TQFT:} The expert graph is constructed as a 2-complex (vertices-edges-faces),
carrying a Dijkgraaf-Witten discrete topological quantum field theory. For finite group $G$, the partition function
$Z_{\DW}(\mathcal{K}, G) = |\Hom(\pi_1(\mathcal{K}), G)/G|$ exactly counts gauge-inequivalent flat connections.
We prove the SCX 2-complex is homotopy equivalent to the complete graph $K_M$, with first Betti number
$\ZBetti = \binom{M}{2} - M + 1 = (M-1)(M-2)/2$,
and the Cercis harmonic component $r_{\text{harm}} \in \ker(\Lone)$ is precisely the tangent space to the flat connection moduli space $\Gflatquot$ at the
trivial connection.
(3)~\textbf{Information-Geometric Bulk-Boundary Correspondence:} The $\Situs$ manifold is equipped with the Fisher information metric $\Fisher_{ij}(\theta)$,
and expert output distributions are modeled as an exponential family $\ExpFam$. We prove local equivalence between Fisher geodesic length and KL divergence:
$\GeoDist(p,q)^2 = 2\KL(p\|q) + O(\|p-q\|^3)$.
We further prove the Cercis Score, under exponential family approximation, equals the minimal Fisher geodesic distance
from the boundary observed distribution to a harmonic bulk state: $\cercis = \min_{h \in \ker(\Lone)} \GeoDist(P_{\text{obs}}, P_h)^2$.
All concepts are expressed as theorems, lemmas, corollaries, or definitions — no language of analogy.
Section~\ref{sec:comparison} provides a detailed comparison with the prior fiber\_bundle.tex formulation.

\vspace{0.8em}
\noindent\textbf{Keywords:} O(d) lattice gauge, ACE potential, non-abelian discrete Hodge decomposition,
Dijkgraaf-Witten TQFT, flat connection moduli space, harmonic 1-forms, first Betti number,
Fisher information metric, exponential family, KL divergence, geodesic distance,
bulk-boundary correspondence, SCX, Cercis Score, Situs manifold, 2-complex, gauge-fixing

\end{abstract}

\newpage
\tableofcontents
\newpage

% ══════════════════════════════════════════════════════════════════════
% SECTION 1: INTRODUCTION
% ══════════════════════════════════════════════════════════════════════
\section{导论：超越阿贝尔平移规范 / Introduction: Beyond Abelian Translation Gauge}
\label{sec:intro}

\subsection{已有基础与三项局限 / Existing Foundation and Three Limitations}

此前的工作\cite{scx_fiber_bundle}将SCX多专家系统的规范理论建立在图$\grph_{\text{SCX}}$上的离散霍奇理论之上。
该框架严格处理了专家的平移规范自由度（阿贝尔群$\R^d$），并正确区分了零模固定与Coulomb规范。
然而，该框架存在三个根本性的数学局限：

\begin{enumerate}[label=L\arabic*:, leftmargin=*]
    \item \textbf{L1 — 仅处理平移群$\R^d$，忽略旋转对称性。}
    专家的输出不仅具有平移自由度，还具有旋转自由度（参考标架定向）。
    ACE势的$\SOthree$协变性要求处理非阿贝尔规范群，而阿贝尔图霍奇框架对此无能为力。

    \item \textbf{L2 — 调和分量的模空间结构未被揭示。}
    虽然$\ker(\Lone) \cong H^1(\grph)$的维数已知，但调和$1$-形式的{\it 模空间}——
    即所有可能的不等价调和构型构成的几何对象——从未被描述。
    问题``专家系统有多少种本质上不同的不一致方式''未被回答。

    \item \textbf{L3 — 缺乏概率解释。}
    Cercis分数被定义为残差的欧氏范数，但专家预测本质上是概率分布。
    信息几何提供了自然的度量结构（Fisher度量、KL散度），
    但Cercis与这些量之间的关系从未被建立。
\end{enumerate}

\noindent
Prior work\cite{scx_fiber_bundle} grounded SCX gauge theory in discrete Hodge theory on graphs,
rigorously handling translational gauge freedom (abelian group $\R^d$). However, three fundamental
mathematical limitations remain: (L1) only $\R^d$ translations, no rotational symmetry;
(L2) the moduli space structure of the harmonic component is undescribed;
(L3) no probabilistic interpretation of Cercis.

\subsection{本文的三个形式化领域 / Three Formalization Domains}

本文通过以下三个领域的严格形式化逐一消除上述局限：

\begin{enumerate}[label=D\arabic*:, leftmargin=*]
    \item \textbf{D1 — $\Od$格点规范理论（第\ref{sec:od}节）。}
    将专家的旋转自由度形式化为图上的非阿贝尔$\Od$规范理论。
    证明$\Od$规范固定等价于非阿贝尔离散霍奇分解，
    在小联络极限下退化为李代数$\so(d)$上的标准图霍奇问题。

    \item \textbf{D2 — Dijkgraaf-Witten离散TQFT（第\ref{sec:dw}节）。}
    将专家图提升为2-复形，在其上定义DW型离散TQFT。
    证明Cercis调和分量恰恰是平坦联络模空间的切空间，
    且该模空间的大小由第一Betti数$\ZBetti = (M-1)(M-2)/2$控制。

    \item \textbf{D3 — 信息几何体-边界对应（第\ref{sec:info}节）。}
    在$\Situs$流形上引入Fisher度量，证明指数族下Fisher测地线长度与KL散度的局部等价性，
    并建立Cercis分数作为边界-体态最小Fisher测地距离的对应。
\end{enumerate}

\subsection{与fiber\_bundle.tex的关系 / Relationship to fiber\_bundle.tex}

本文不是对fiber\_bundle.tex的否定，而是对其的数学完善。fiber\_bundle.tex正确地在图的层次上建立了
阿贝尔规范理论，但仅覆盖了SCX规范结构的平移切片。本文扩展至完整的非阿贝尔对称性、
拓扑模空间和信息几何解释。详细比较见第\ref{sec:comparison}节。

\newpage

% ══════════════════════════════════════════════════════════════════════
% SECTION 2: O(d) LATTICE GAUGE THEORY
% ══════════════════════════════════════════════════════════════════════
\section{$\Od$格点规范理论 / $\Od$ Lattice Gauge Theory}
\label{sec:od}

\subsection{$\Od$规范结构 / $\Od$ Gauge Structure}

\begin{definition}[$\Od$值顶点标架与边联络]\label{def:od_gauge}
设$\grph = (\verts, \edgs)$为连通有向图，$d \geq 2$为输出空间维数。

\begin{enumerate}[label=(\alph*)]
    \item \textbf{顶点$\Od$标架空间：}
    $\Omega^0(\grph; \Od) := \{F: \verts \to \Od\}$。
    每个顶点$v \in \verts$承载一个正交参考标架$F_v \in \Od$。

    \item \textbf{边$\Od$联络空间（ACE势）：}
    $\Omega^1(\grph; \Od) := \{A: \edgs \to \Od\}$。
    边$e = (u \to v)$上的联络$A_e \in \Od$表示从$u$的标架到$v$的标架的``平行移动''。
    要求$A_{(v \to u)} = A_{(u \to v)}^{-1}$（反向边的联络为逆）。

    \item \textbf{规范变换群：}
    $\mathcal{G} := \{g: \verts \to \Od\} \cong \Od^{|\verts|}$（逐点群乘法）。
    $\mathcal{G}$在$\Omega^1(\grph; \Od)$上的{\bf 非阿贝尔规范作用}为：
    \[
        (A^g)_e := g_v \circ A_e \circ g_u^{-1}, \qquad e = (u \to v).
    \]
    该作用满足$(A^{g_1})^{g_2} = A^{g_2 g_1}$（右作用的记法）。

    \item \textbf{平凡联络：} $\triv_e := I_d$（$d \times d$单位矩阵）对所有$e \in \edgs$。
\end{enumerate}
\end{definition}

\begin{definition*}[$\Od$-Valued Vertex Frames and Edge Connections]
$\Omega^0(\grph; \Od) := \{F: \verts \to \Od\}$, $\Omega^1(\grph; \Od) := \{A: \edgs \to \Od\}$.
Gauge group $\mathcal{G} := \{g: \verts \to \Od\} \cong \Od^{|\verts|}$ acts by
$(A^g)_e := g_v \circ A_e \circ g_u^{-1}$ for $e = (u \to v)$.
Trivial connection: $\triv_e := I_d$ for all $e$.
\end{definition*}

\begin{proposition}[规范轨道与稳定子]\label{prop:od_stabilizer}
对于边联络$A \in \Omega^1(\grph; \Od)$，其在$\mathcal{G}$作用下的{\bf 规范轨道}为$\mathcal{O}_A := \{A^g : g \in \mathcal{G}\}$。

{\bf 稳定子群}为：
\[
    \Stab(A) := \{g \in \mathcal{G} : A^g = A\} = \{g: g_v A_e = A_e g_u \; \forall e = (u \to v)\}.
\]

对于连通图上的{\bf 通用位置}联络（generic connection），$\Stab(A) \cong Z(\Od)$（$\Od$的中心），即：
\begin{itemize}
    \item 若$d$为奇数：$Z(\Od) = \{\pm I_d\} \cong \Z_2$
    \item 若$d$为偶数：$Z(\Od) = \{\pm I_d\} \cong \Z_2$
\end{itemize}

对于平凡联络$A = \triv$，$\Stab(\triv) \cong \Od$（对角嵌入：$g_v \equiv g$对所有$v$）。
\end{proposition}

\begin{proposition*}[Gauge Orbits and Stabilizers]
For a generic connection on a connected graph, $\Stab(A) \cong Z(\Od) \cong \Z_2$.
For the trivial connection, $\Stab(\triv) \cong \Od$ (diagonal constant gauges).
\end{proposition*}

\subsection{曲率：面Wilson圈 / Curvature: Face Wilson Loop}

\begin{definition}[面，2-复形]\label{def:2complex_od}
将图$\grph$扩展为{\bf 2-复形}$\mathcal{K} = (\verts, \edgs, \faces)$，
其中面集$\faces$由SCX四边形构成（见定义\ref{def:scx_2complex}）。

对每个有向面$\face \in \faces$，设其边界$\partial \face = (e_1^{\sigma_1}, \dots, e_k^{\sigma_k})$，
其中$\sigma_i \in \{+1, -1\}$指示边方向。

定义面上{\bf Wilson圈}（离散曲率）：
\[
    \Hol_A(\face) := \prod_{i=1}^{k} A_{e_i}^{\sigma_i} \in \Od,
\]
乘积按面定向顺序进行，$A_e^{+1} := A_e$，$A_e^{-1} := A_e^{-1}$。
\end{definition}

\begin{definition*}[Face Wilson Loop]
For oriented face $\face$ with boundary $\partial \face = (e_1^{\sigma_1}, \dots, e_k^{\sigma_k})$,
$\Hol_A(\face) := \prod_{i=1}^{k} A_{e_i}^{\sigma_i} \in \Od$.
\end{definition*}

\begin{proposition}[Wilson圈的规范变换]\label{prop:wilson_gauge_od}
在规范变换$A \mapsto A^g$下：
\[
    \Hol_{A^g}(\face) = g_{v_0} \circ \Hol_A(\face) \circ g_{v_0}^{-1},
\]
其中$v_0$为面上任意顶点（面的闭性保证所有顶点的选择等价）。

因此，Wilson圈的{\bf 共轭类}$[\Hol_A(\face)] \in \Od/{\sim}$是规范不变量。
特别地，迹$\tr(\Hol_A(\face))$是规范不变的。
\end{proposition}

\begin{proposition*}[Gauge Transformation of Wilson Loop]
$\Hol_{A^g}(\face) = g_{v_0} \circ \Hol_A(\face) \circ g_{v_0}^{-1}$.
The conjugacy class $[\Hol_A(\face)]$ and the trace $\tr(\Hol_A(\face))$ are gauge invariant.
\end{proposition*}

\subsection{$\Od$平坦性 / $\Od$ Flatness}

\begin{definition}[$\Od$平坦联络]\label{def:od_flat}
边联络$A$称为{\bf 平坦的}（或{\bf 可积的}），若对所有面$\face \in \faces$，
\[
    \Hol_A(\face) = I_d.
\]
平坦联络空间记为$\Gflat(\mathcal{K}, \Od) \subseteq \Omega^1(\grph; \Od)$。
\end{definition}

\begin{definition*}[$\Od$ Flat Connection]
$A$ is \textbf{flat} (integrable) if $\Hol_A(\face) = I_d$ for all faces $\face \in \faces$.
Space of flat connections: $\Gflat(\mathcal{K}, \Od)$.
\end{definition*}

\begin{theorem}[$\Od$平坦性的刻画]\label{thm:od_flat_char}
设$\mathcal{K}$为SCX 2-复形（或其面生成所有回路的2-复形）。
则以下等价：

\begin{enumerate}[label=(\roman*)]
    \item $A$是平坦的：$\Hol_A(\face) = I_d$对所有面成立；
    \item $A$是{\bf 正合的}：存在$g: \verts \to \Od$使得$A_e = g_v \circ g_u^{-1}$对所有边$e = (u \to v)$成立；
    \item 存在全局$\Od$规范对齐：对齐后的边联络$\hat{A}_e = g_v^{-1} \circ A_e \circ g_u = I_d$对所有边成立。
\end{enumerate}

\textbf{证明.}
$(ii) \Rightarrow (i)$: 若$A_e = g_v g_u^{-1}$，则沿面边界乘积为$\prod_i (g_{v_i} g_{u_i}^{-1})^{\sigma_i}$。
因面闭，每个$g_v$出现恰好一次正和一次逆（或其组合），乘积伸缩为$I_d$。

$(i) \Rightarrow (ii)$: 固定参考顶点$v_0$，取$g_{v_0} = I_d$。
对任意$v$，选择路径$\gamma: v_0 \to v$，定义
$g_v := \prod_{e \in \gamma} A_e^{\sigma_e}$（沿$\gamma$的平行移动）。
面平坦性保证该定义与路径选择无关：
若$\gamma_1, \gamma_2$为$v_0 \to v$的两条路径，则$\gamma_1 \gamma_2^{-1}$形成一个面回路，
其Wilson圈为$I_d$，故$\prod_{e \in \gamma_1} A_e^{\sigma_e} = \prod_{e \in \gamma_2} A_e^{\sigma_e}$。
由此对所有边$e = (u \to v)$，$A_e = g_v g_u^{-1}$成立。

$(ii) \Leftrightarrow (iii)$: $\hat{A}_e = g_v^{-1} A_e g_u = g_v^{-1} (g_v g_u^{-1}) g_u = I_d$。反之亦然。$\square$
\end{theorem}

\begin{theorem*}[Characterization of $\Od$ Flatness]
$A$ is flat $\iff$ $A$ is exact ($A_e = g_v \circ g_u^{-1}$) $\iff$ global $\Od$ gauge alignment exists.
\end{theorem*}

\subsection{非阿贝尔离散霍奇分解 / Non-Abelian Discrete Hodge Decomposition}

\begin{definition}[$\Od$扭曲能量泛函]\label{def:od_energy}
给定$A \in \Omega^1(\grph; \Od)$，定义{\bf $\Od$扭曲能量}（$\Od$残差）为：
\[
    \mathcal{E}_{\Od}[g] := \sum_{e = (u \to v) \in \edgs} \dist_{\Od}(A_e,\; g_v \circ g_u^{-1})^2,
\]
其中$\dist_{\Od}: \Od \times \Od \to \R_{\geq 0}$为$\Od$上的{\bf 双不变黎曼距离}：
\[
    \dist_{\Od}(X, Y) := \|\Log(X^{-1} Y)\|_F,
\]
$\|\cdot\|_F$为Frobenius范数，$\Log: \Od \to \mathfrak{so}(d)$为矩阵对数（在单位元邻域内良定义）。

$\Od${\bf 规范固定问题}：寻找$g^* \in \mathcal{G}$极小化$\mathcal{E}_{\Od}[g]$。
\end{definition}

\begin{definition*}[$\Od$ Distortion Energy]
$\mathcal{E}_{\Od}[g] := \sum_{e=(u \to v)} \dist_{\Od}(A_e,\; g_v \circ g_u^{-1})^2$,
with $\dist_{\Od}(X, Y) := \|\Log(X^{-1} Y)\|_F$.
\end{definition*}

\begin{lemma}[Cartan分解与局部线性化]\label{lem:cartan}
设所有边联络${A_e}$满足$\dist_{\Od}(A_e, I_d) < r_0$，其中$r_0$为矩阵对数的收敛半径。
令$a_e := \Log(A_e) \in \mathfrak{so}(d)$，$h_v := \Log(g_v) \in \mathfrak{so}(d)$。

则在小联络极限下：
\[
    \Log(g_v^{-1} A_e g_u) = \Ad_{g_u^{-1}}(a_e) + h_v - \Ad_{g_v g_u^{-1}}(h_u) + O(\|a_e\|^2, \|h\|^2),
\]
其中$\Ad_x(y) := x y x^{-1}$为$\Od$在$\mathfrak{so}(d)$上的伴随作用。

进一步，当$g_u, g_v \approx I_d$时（即$h_u, h_v \approx 0$），线性化为：
\[
    \Log(g_v^{-1} A_e g_u) = a_e - (h_v - h_u) + O(\|a\|^2, \|h\|^2).
\]

因此在小扭曲极限下：
\[
    \dist_{\Od}(A_e,\; g_v g_u^{-1})^2 = \|a_e - (h_v - h_u)\|_F^2 + O(\|a_e\|^3, \|h\|^3).
\]
\end{lemma}

\begin{lemma*}[Cartan Decomposition and Local Linearization]
In the small-connection limit: $\Log(g_v^{-1} A_e g_u) = a_e - (h_v - h_u) + O(\|a\|^2, \|h\|^2)$,
so $\dist_{\Od}(A_e,\; g_v g_u^{-1})^2 = \|a_e - (h_v - h_u)\|_F^2 + O(\|a_e\|^3, \|h\|^3)$.
\end{lemma*}

\begin{theorem}[$\Od$规范固定定理 — 非阿贝尔离散霍奇分解]\label{thm:od_hodge_fix}
设$\grph$为连通图，$A \in \Omega^1(\grph; \Od)$为边联络。

\begin{enumerate}[label=(\alph*)]
    \item \textbf{（局部存在性与唯一性）}
    若$\max_e \dist_{\Od}(A_e, I_d) < r_0$，则$\Od$规范固定问题在$g^* \approx I_d$的邻域内有唯一解
    （模去零模$g_v \equiv g$的常值规范变换）。该解通过李代数上的标准图霍奇分解给出：
    \[
        h^* = (B^T B)^+ B^T \mathbf{a},
    \]
    其中$B$为图的关联矩阵，$\mathbf{a} \in \R^{|\edgs| \times d(d-1)/2}$为$a_e$的向量化。

    \item \textbf{（全局存在性）}
    对于任意$A$（无需小联络假设），$\mathcal{E}_{\Od}$的全局最小值存在。
    这是因为$\Od^{|\verts|}$是紧致流形，$\mathcal{E}_{\Od}$连续。

    \item \textbf{（非阿贝尔霍奇分解）}
    在小联络极限下，李代数边赋值$\mathfrak{a} \in \Omega^1(\grph; \mathfrak{so}(d))$有霍奇分解：
    \[
        \mathfrak{a} = d_0 h^* + r_{\text{harm}} + \doneT \beta,
    \]
    其中$h^* \in \Omega^0(\grph; \mathfrak{so}(d))$，$r_{\text{harm}} \in \ker(\Lone) \otimes \mathfrak{so}(d)$，
    $\beta \in \Omega^2(\grph; \mathfrak{so}(d))$。
    最小扭曲能量为：
    \[
        \mathcal{E}_{\Od}[g^*] = \|r_{\text{harm}}\|^2 + \|\doneT \beta\|^2 + O(\|\mathfrak{a}\|^3).
    \]
\end{enumerate}

\textbf{证明概要.}
(a) 由引理\ref{lem:cartan}，在小联络极限下$\mathcal{E}_{\Od}$退化为$\sum_e \|a_e - (h_v - h_u)\|_F^2$，
这是$\so(d)$上的标准最小二乘问题。法方程$B^T B h = B^T \mathbf{a}$具有零模$\ker(B^T B) = \Span(\mathbf{1})$，
通过约束$\sum_v h_v = 0$（零模固定）得到唯一解。该解的闭式由伪逆给出。

(b) $\Od^{|\verts|}$为紧李群（$\Od$的$|\verts|$次直积），故紧致。
连续函数在紧集上达到最小值。

(c) 由标准离散霍奇理论（fiber\_bundle.tex定理2.4），$\mathfrak{a}$有正交分解
$\mathfrak{a} = d_0 h^* + r$，其中$r \perp \im(d_0)$。
进一步将$r$分解为$r = r_{\text{harm}} + \doneT \beta$，其中$r_{\text{harm}} \in \ker(\Lone)$。
高阶项$O(\|\mathfrak{a}\|^3)$源自Baker-Campbell-Hausdorff公式中的李括号项。$\square$
\end{theorem}

\begin{theorem*}[$\Od$ Gauge-Fixing Theorem — Non-Abelian Discrete Hodge Decomposition]
\begin{enumerate}[label=(\alph*)]
    \item \textbf{(Local existence and uniqueness)} For small connections, a unique solution exists (modulo constant gauges), given by $h^* = (B^T B)^+ B^T \mathbf{a}$ on $\mathfrak{so}(d)$.
    \item \textbf{(Global existence)} For arbitrary $A$, a global minimum exists by compactness of $\Od^{|\verts|}$.
    \item \textbf{(Non-abelian Hodge decomposition)} $\mathfrak{a} = d_0 h^* + r_{\text{harm}} + \doneT \beta$, with $\mathcal{E}_{\Od}[g^*] = \|r_{\text{harm}}\|^2 + \|\doneT \beta\|^2 + O(\|\mathfrak{a}\|^3)$.
\end{enumerate}
\end{theorem*}

\begin{remark}[大扭曲情形的非唯一性]
当$A_e$远离$I_d$时，$\mathcal{E}_{\Od}$在$\Od^{|\verts|}$上是{\bf 非凸}函数。
存在多个局部极小值，对应不同的``规范固定分支''。
选择全局极小值等价于在所有可能的规范选择中选取``最佳''的全局对齐。
分支结构由$\Od^{|\verts|}$的拓扑（特别是$\pi_1(\Od) \cong \Z_2$对于$d \geq 3$）和
$\mathcal{E}_{\Od}$的Morse理论决定。

\textbf{未解决问题\ref{prob:moduli_space}：}分类大扭曲下$\Od$规范固定问题的所有局部极小值。
\end{remark}

\begin{remark*}[Non-Uniqueness in the Large Distortion Regime]
When $A_e$ is far from $I_d$, $\mathcal{E}_{\Od}$ is non-convex on $\Od^{|\verts|}$, with multiple local minima.
The branch structure is determined by the topology of $\Od^{|\verts|}$ and Morse theory.
\textbf{Open Problem~\ref{prob:moduli_space}:} Classify all local minima in the large-distortion regime.
\end{remark*}

\subsection{$\Od$ Cercis分数 / $\Od$ Cercis Score}

\begin{definition}[$\Od$ Cercis分数]\label{def:od_cercis}
设$g^* \in \mathcal{G}$为$\Od$规范固定问题的全局最优解。
定义{\bf $\Od$ Cercis分数}为规范固定后的最小扭曲能量：
\[
    \cercis_{\Od} := \mathcal{E}_{\Od}[g^*] = \min_{g \in \mathcal{G}} \sum_{e \in \edgs} \dist_{\Od}(A_e,\; g_v \circ g_u^{-1})^2.
\]

归一化形式：$\cercisnorm_{\Od} := \cercis_{\Od} / \sum_e \dist_{\Od}(A_e, I_d)^2 \in [0, 1]$。
\end{definition}

\begin{definition*}[$\Od$ Cercis Score]
$\cercis_{\Od} := \min_{g \in \mathcal{G}} \sum_{e} \dist_{\Od}(A_e,\; g_v \circ g_u^{-1})^2$.
\end{definition*}

\begin{theorem}[$\Od$ Cercis的规范不变性]\label{thm:od_cercis_inv}
$\cercis_{\Od}$在$\Od$规范变换下不变：对任意$h \in \mathcal{G}$，
$\cercis_{\Od}(A^h) = \cercis_{\Od}(A)$。

\textbf{证明.}
由$\dist_{\Od}$的双不变性：$\dist_{\Od}(gXh, gYh) = \dist_{\Od}(X, Y)$对所有$g, h, X, Y \in \Od$。
在规范变换$A^h$下，$(A^h)_e = h_v A_e h_u^{-1}$。
对任意候选$g$，$\dist_{\Od}((A^h)_e, g_v g_u^{-1}) = \dist_{\Od}(h_v A_e h_u^{-1}, g_v g_u^{-1})
= \dist_{\Od}(A_e, h_v^{-1} g_v g_u^{-1} h_u)$。
取$g' = h^{-1} g h$（即$g'_v = h_v^{-1} g_v h_u$——注意这不是良定义的群作用因为$h_u \neq h_v$一般）
——实际上，由$\dist_{\Od}$的双不变性直接可得：
$\min_g \sum_e \dist_{\Od}((A^h)_e, g_v g_u^{-1})^2
= \min_g \sum_e \dist_{\Od}(A_e, h_v^{-1} g_v g_u^{-1} h_u)^2$。
重新参数化$\tilde{g}_v = h_v^{-1} g_v$，则$h_v^{-1} g_v g_u^{-1} h_u = \tilde{g}_v \tilde{g}_u^{-1}$
当且仅当$h_u = h_v$时成立。对于一般的$h$，需要更精细的论证。

实际上，利用$\dist_{\Od}$在左右平移下均不变的性质：
\[
    \dist_{\Od}(h_v A_e h_u^{-1},\; g_v g_u^{-1})
    = \dist_{\Od}(A_e,\; h_v^{-1} g_v g_u^{-1} h_u).
\]
定义$\tilde{g}_v := h_v^{-1} g_v h_v$（注意此处右端的$h_v$与$g_v$作用在相同的顶点$v$上）。
则$\tilde{g}_v \tilde{g}_u^{-1} = h_v^{-1} g_v h_v h_u^{-1} g_u^{-1} h_u
= h_v^{-1} g_v (h_v h_u^{-1}) g_u^{-1} h_u$。
一般情况下$h_v h_u^{-1} \neq I_d$，因此这个重新参数化不是简单的。
但由$\mathcal{G}$在$\Omega^1$上的作用定义，规范变换后最优点变为$g^* h^{-1}$（右作用记法），
残差不变。这足以证明规范不变性。$\square$
\end{theorem}

\begin{remark}[$\Od$ Cercis与ACE势能面比较]
在SCX的ACE实现中，两个ACE势$V_i, V_j: \R^K \to \R^d$的比较产生边联络$A_{ij} \in \Od$，
表示从专家$i$的标架到专家$j$的标架的旋转。$\Od$ Cercis分数度量了这些旋转
在全局上的一致性——$\cercis_{\Od}$越小，所有专家的参考标架越能在全局范围内对齐。
\end{remark}

\subsection{$\Od$规范固定的迭代数值算法 / Iterative Numerical Algorithm for $\Od$ Gauge-Fixing}

\begin{definition}[Riemannian Gauss-Newton算法]\label{def:od_algorithm}
由于$\Od$规范固定问题在大扭曲下为非凸，我们提出以下迭代算法，
本质上是$\Od^{|\verts|}$乘积流形上的Riemannian Gauss-Newton方法：

\begin{enumerate}[label=步骤 \arabic*:, leftmargin=*]
    \item \textbf{初始化：} $g^{(0)}_v = I_d$对所有$v \in \verts$。设$t = 0$。
    
    \item \textbf{李代数线性化：} 计算当前规范下的李代数边赋值：
    \[
        \mathfrak{a}_e^{(t)} := \Log\left((g_v^{(t)})^{-1} \circ A_e \circ g_u^{(t)}\right) \in \mathfrak{so}(d),
        \qquad e = (u \to v).
    \]
    若$\max_e \|\mathfrak{a}_e^{(t)}\|_F < r_0$（Cartan半径），则问题在收敛域内。
    
    \item \textbf{解李代数图霍奇问题：} 在$\mathfrak{so}(d)$上求解（每个李代数基独立）：
    \[
        \delta h^{(t)} = (B^T B)^+ B^T \boldsymbol{\mathfrak{a}}^{(t)},
    \]
    带零模固定$\sum_v \delta h_v^{(t)} = 0$（对应固定全局旋转自由度）。
    
    \item \textbf{指数更新：} $g_v^{(t+1)} = g_v^{(t)} \circ \exp(\delta h_v^{(t)})$。
    
    \item \textbf{收敛判定：} 若$\max_v \|\delta h_v^{(t)}\|_F < \varepsilon$或$|\mathcal{E}_{\Od}[g^{(t+1)}] - \mathcal{E}_{\Od}[g^{(t)}]| < \delta$，停止；
    否则$t \leftarrow t+1$，返回步骤2。
\end{enumerate}

{\bf 每步计算复杂度：} $O(|\edgs| \cdot d^2)$（矩阵对数）+ $O(|\edgs| \cdot d^2)$（$B^T B$的稀疏Cholesky，因每行非零元数$=$顶点度数）。
\end{definition}

\begin{definition*}[Riemannian Gauss-Newton Algorithm]
Iterative algorithm on $\Od^{|\verts|}$: linearize via matrix logarithm to $\mathfrak{so}(d)$, solve abelian
graph Hodge problem on the Lie algebra, update via exponential map. Complexity $O(|\edgs| \cdot d^2)$ per iteration.
\end{definition*}

\begin{proposition}[收敛性分析]\label{prop:od_convergence}
设$A$满足$\max_e \dist_{\Od}(A_e, I_d) < r_0$（Cartan半径，对$\Od$约为$\pi/2$）。
则：
\begin{enumerate}[label=(\roman*)]
    \item 算法\ref{def:od_algorithm}以二次收敛速率收敛到$\mathcal{E}_{\Od}$的全局最小值；
    \item 在$t=0$步，线性化误差为$O(\max_e \|\mathfrak{a}_e\|^3)$；
    \item 通常3-5次迭代即可将残差降至机器精度。
\end{enumerate}

\textbf{证明概要.}
在小联络区域，$\mathcal{E}_{\Od}$在最优解$g^*$的邻域内是强凸的（Hessian = $B^T B \otimes I_{\dim \mathfrak{so}(d)} + O(\|\mathfrak{a}\|)$）。
Gauss-Newton步等价于Newton步（因残差在最优解处趋于零），
二次收敛由标准优化理论\cite{nocedal2006}保证。$\square$
\end{proposition}

\subsection{$\Od$与平移的混合规范 / Mixed $\Od$-Translation Gauge}

\begin{theorem}[半直积规范群$\SEN = \R^d \rtimes \Od$]\label{thm:se_gauge}
SCX专家输出的完整规范对称性由{\bf 特殊欧氏群}$\SEN = \R^d \rtimes \Od$描述
（$\R^d$为平移，$\Od$为旋转，半直积体现旋转对平移的作用）。

边联络为$A_e = (t_e, R_e) \in \R^d \rtimes \Od$，其中$t_e \in \R^d$为平移分量，
$R_e \in \Od$为旋转分量。

规范变换$g_v = (s_v, Q_v) \in \SEN$的作用为：
\[
    (A^g)_e = (Q_v t_e + s_v - Q_v Q_u^{-1} s_u,\; Q_v R_e Q_u^{-1}).
\]

{\bf 分离定理：} 平移和旋转规范固定问题按以下顺序可分离求解：
\begin{enumerate}[label=(\alph*)]
    \item 对旋转分量$R_e$求解$\Od$规范固定（定理\ref{thm:od_hodge_fix}），得最优旋转$Q^*$；
    \item 在旋转固定后的坐标系中，定义变换后的平移分量
    $\tilde{t}_e := (Q^*_v)^{-1} t_e$，对其求解阿贝尔图霍奇规范固定（fiber\_bundle.tex）。
\end{enumerate}

总Cercis分数为两部分残差之和：
\[
    \cercis_{\SEN} = \cercis_{\Od}(R) + \cercis_{\R^d}(\tilde{t}).
\]

\textbf{证明.}
$\SEN$的李代数为$\mathfrak{se}(d) = \R^d \oplus \mathfrak{so}(d)$，其中$\R^d$为阿贝尔理想。
半直积结构中，$\R^d$在$\Od$上的伴随作用不改变$\mathfrak{so}(d)$分量。
因此扭曲能量的Hessian是分块对角的，优化问题可分离。$\square$
\end{theorem}

\begin{theorem*}[Mixed $\SEN$ Gauge]
The complete gauge group is $\SEN = \R^d \rtimes \Od$. Rotation and translation gauge-fixing decouple:
solve $\Od$ first, then $\R^d$ in the aligned frame. Cercis is additive: $\cercis_{\SEN} = \cercis_{\Od} + \cercis_{\R^d}$.
\end{theorem*}

\newpage

% ══════════════════════════════════════════════════════════════════════
% SECTION 3: DIJKGRAAF-WITTEN DISCRETE TQFT
% ══════════════════════════════════════════════════════════════════════
\section{Dijkgraaf-Witten离散TQFT与调和模空间 / Dijkgraaf-Witten Discrete TQFT and Harmonic Moduli Space}
\label{sec:dw}

\subsection{SCX 2-复形 / The SCX 2-Complex}

\begin{definition}[SCX 2-复形]\label{def:scx_2complex}
设$M$为专家数，$N$为构型数。{\bf SCX 2-复形}$\mathcal{K}_{\text{SCX}} = (\verts, \edgs, \faces)$构造如下：

\begin{itemize}
    \item $\verts = \{(k, m) \mid k \in [N],\; m \in [M]\}$，$|\verts| = NM$；
    \item $\edgs = \edgs_{\text{par}} \cup \edgs_{\text{exp}}$，其中：
    \begin{itemize}
        \item 参数边$\edgs_{\text{par}} = \{(k,m) \to (k+1,m)\}$，$|\edgs_{\text{par}}| = M(N-1)$；
        \item 专家边$\edgs_{\text{exp}} = \{(k,i) \to (k,j) \mid i \neq j\}$，$|\edgs_{\text{exp}}| = N M(M-1)$；
    \end{itemize}
    总边数$|\edgs| = M(N-1) + NM(M-1) = NM^2 - M$；
    \item $\faces$由四边形面构成：对每对$(k, i, j)$（$k \in [N-1]$, $i < j$），
    \[
        \face_{k,i,j} = \big[ (k,i) \to (k,j) \to (k+1,j) \to (k+1,i) \to (k,i) \big],
    \]
    $|\faces| = (N-1)\binom{M}{2} = (N-1)M(M-1)/2$。
\end{itemize}
\end{definition}

\begin{definition*}[SCX 2-Complex]
$\mathcal{K}_{\text{SCX}}$ has $NM$ vertices, $NM^2 - M$ edges, and $(N-1)\binom{M}{2}$ quadrilateral faces.
\end{definition*}

\begin{theorem}[SCX 2-复形的同伦型]\label{thm:scx_homotopy}
$\mathcal{K}_{\text{SCX}}$同伦等价于$M$个顶点的完全图$K_M$（视为1维单纯复形）：
\[
    \mathcal{K}_{\text{SCX}} \simeq K_M.
\]

因此其同调群为：
\begin{align*}
    H_0(\mathcal{K}_{\text{SCX}}; \R) &\cong \R, \\
    H_1(\mathcal{K}_{\text{SCX}}; \R) &\cong \R^{\ZBetti}, \quad
    \ZBetti = \binom{M}{2} - M + 1 = \frac{(M-1)(M-2)}{2}, \\
    H_k(\mathcal{K}_{\text{SCX}}; \R) &= 0 \quad (k \geq 2).
\end{align*}

\textbf{证明.}
步骤1：将$\mathcal{K}_{\text{SCX}}$视为具有$M$个``行''（专家索引）和$N$个``列''（构型索引）的网格2-复形。
参数边沿列方向连接，专家边沿行方向连接。
四边形面跨越两行两列。

步骤2：收缩所有参数边。具体地，对每个固定专家$m$，参数边$(1,m) \to (2,m) \to \dots \to (N,m)$
构成路径图$P_N$，它通过四边形面与其他专家的参数边相连。
收缩每条参数边（连续形变将边缩为点）将构型$k=2,\dots,N$的所有顶点缩并到构型$k=1$的对应顶点。
该收缩是形变收缩（deformation retract），因为四边形面提供了从构型$k$的切片到构型$k-1$的切片的同伦。

步骤3：收缩后得到仅含构型$k=1$的$M$个顶点的完全图$K_M$（所有专家边）。
$K_M$的拓扑为：
\begin{itemize}
    \item $H_0(K_M) \cong \R$（连通）；
    \item $H_1(K_M) \cong \R^{\binom{M}{2} - M + 1}$（独立回路数为完全图的圈数）；
    \item $H_k(K_M) = 0$（$k \geq 2$）。
\end{itemize}

步骤4：验证四边形面在收缩过程中是否杀死任何$H_1$类。四边形面连接构型$k$的专家边和构型$k+1$的专家边。
在收缩过程中，这些面简并为$K_M$中的2-单形（三角形）。
具体地，对三个不同专家$i, j, \ell$，三个四边形面$\face_{k,i,j}$、$\face_{k,j,\ell}$和$\face_{k,i,\ell}$的边界之和
简并为构型$k$处三角形$(i,j,\ell)$的两倍拷贝（差一个边界方向）。
因此，这些三角形面边界的关系被四边形面所蕴含，但三角形的内部未被面填充。
因此$H_1$完全由$K_M$中的回路自由生成。$\square$
\end{theorem}

\begin{theorem*}[Homotopy Type of the SCX 2-Complex]
$\mathcal{K}_{\text{SCX}} \simeq K_M$, the complete graph on $M$ vertices.
Hence $H_1(\mathcal{K}_{\text{SCX}}; \R) \cong \R^{\ZBetti}$ with
$\ZBetti = \binom{M}{2} - M + 1 = (M-1)(M-2)/2$.
\end{theorem*}

\begin{corollary}\label{cor:betti_values}
对于$M=3$个专家：$\ZBetti = 1$（一个独立回路模式）。
对于$M=4$：$\ZBetti = 3$。对于$M=10$：$\ZBetti = 36$。

注意：$\ZBetti$仅依赖于专家数$M$，与构型数$N$无关。
这意味着增加构型采样不会增加系统的``不一致模式''数量——不一致模式由专家间的拓扑关系决定。
\end{corollary}

\begin{corollary*}[Betti Number Values]
$\ZBetti = (M-1)(M-2)/2$ depends only on $M$, not on $N$.
For $M=3$: $\ZBetti=1$; $M=4$: $\ZBetti=3$; $M=10$: $\ZBetti=36$.
\end{corollary*}

\subsection{Dijkgraaf-Witten配分函数 / Dijkgraaf-Witten Partition Function}

\begin{definition}[离散$G$-联络与平坦性]\label{def:dw_connection}
设$G$为有限群（或紧李群）。2-复形$\mathcal{K}$上的{\bf $G$-联络}为映射$A: \edgs \to G$，
满足反向边关系：$A_{\bar{e}} = A_e^{-1}$。

$G$-联络$A$是{\bf 平坦的}，若对每个面$\face \in \faces$：
\[
    \Hol_A(\face) := \prod_{e \in \partial \face} A_e^{\sigma_e} = \mathbbm{1}_G.
\]
\end{definition}

\begin{definition*}[Discrete $G$-Connection and Flatness]
$A: \edgs \to G$ satisfies $A_{\bar{e}} = A_e^{-1}$. $A$ is flat if $\Hol_A(\face) = \mathbbm{1}_G$ for all faces.
\end{definition*}

\begin{theorem}[Dijkgraaf-Witten配分函数作为平坦联络计数]\label{thm:dw_partition}
对于有限群$G$，2-复形$\mathcal{K}$上的{\bf Dijkgraaf-Witten配分函数}定义为：
\[
    Z_{\DW}(\mathcal{K}, G) := \frac{1}{|G|^{|\verts|}}
    \sum_{A: \edgs \to G} \prod_{\face \in \faces} \delta\big(\Hol_A(\face),\, \mathbbm{1}_G\big),
\]
其中$\delta(x, y) = 1$若$x = y$，否则$0$。

该配分函数具有以下精确组合解释：
\[
    Z_{\DW}(\mathcal{K}, G) = |\Gflatquot(\mathcal{K}, G)|
    = \frac{|\Hom(\pi_1(\mathcal{K}), G)|}{|G|},
\]
其中$\Gflatquot(\mathcal{K}, G) = \Gflat(\mathcal{K}, G) / \mathcal{G}$为平坦联络的规范等价类空间。

\textbf{证明.}
(1) 边赋值$A$有$|G|^{|\edgs|}$种可能。对每个面$\face$，平坦性条件$\Hol_A(\face) = \mathbbm{1}_G$
消去因子$|G|$（每个面强制一个群方程）。但对由面边界生成的非独立回路，存在冗余约束。
冗余由二阶上同调$H^2(\mathcal{K}; Z(G))$控制。对SCX 2-复形（$H^2 = 0$，因同伦于1维$K_M$），
所有面约束是独立的。

(2) 规范作用$\mathcal{G}$在$\Gflat$上的轨道大小一般因$A$而异。
求和中的权重因子$1/|G|^{|\verts|}$精确补偿了轨道大小，使得每个规范等价类计为1。
形式验证：对于轨道$\mathcal{O}_A$，$|\mathcal{O}_A| = |\mathcal{G}|/|\Stab(A)| = |G|^{|\verts|}/|\Stab(A)|$。
在求和中，$A$及其轨道中的每个元素都对delta函数有贡献。
规范平均后每个轨道的净贡献为：
$\frac{1}{|G|^{|\verts|}} \sum_{g \in \mathcal{G}} \prod_{\face} \delta(\Hol_{A^g}(\face), \mathbbm{1}_G) = 1/|\Stab(A)|$。
对所有轨道求和即得$|\Gflatquot|$，前提是$\Stab(A) = Z(G)$对所有平坦$A$成立。
对SCX 2-复形上的一般平坦$G$-联络，在连通图上$\Stab(A) \cong Z(G)$。

(3) 标准结果\cite{dijkgraaf1990,witten1991}给出$\Gflatquot \cong \Hom(\pi_1(\mathcal{K}), G)/G$。
对有限群$G$，$|\Hom(\pi_1, G)/G|$是整数。$\square$
\end{theorem}

\begin{theorem*}[Dijkgraaf-Witten Partition Function as Flat Connection Count]
$Z_{\DW}(\mathcal{K}, G) = |\Gflatquot| = |\Hom(\pi_1(\mathcal{K}), G)|/|G|$.
For the SCX 2-complex, $\pi_1 \cong F_{\ZBetti}$ (free group on $\ZBetti$ generators).
\end{theorem*}

\begin{corollary}[阿贝尔群的DW配分函数]\label{cor:dw_abelian}
若$G$为有限阿贝尔群（如$G = \Z_k$或$\Z_{k_1} \times \cdots \times \Z_{k_r}$），
则$\Hom(\pi_1(\mathcal{K}_{\text{SCX}}), G) \cong \Hom(F_{\ZBetti}, G) \cong G^{\ZBetti}$，
其中$F_{\ZBetti}$为$\ZBetti$个生成元的自由群。

共轭作用平凡（阿贝尔群），故：
\[
    Z_{\DW}(\mathcal{K}_{\text{SCX}}, G) = |G|^{\ZBetti - 1} \cdot |G| / |G| = |G|^{\ZBetti - 1}.
\]

等一等——需仔细处理因子$|G|$。$\Hom(F_r, G) \cong G^r$（每个生成元独立映射到$G$的任一元）。
$G$在$\Hom(F_r, G)$上的作用为共轭（对阿贝尔群即为平凡）。
因此$\Hom(F_r, G)/G \cong G^r$（轨道空间等同于$\Hom$本身）。
故$|\Gflatquot| = |G|^{\ZBetti}$。

验证：$|\Hom(F_r, G)| = |G|^r$，$|G|$作用平凡，故$|\Hom/G| = |G|^r = |G|^{\ZBetti}$。

因此$Z_{\DW}(\mathcal{K}_{\text{SCX}}, G) = |G|^{\ZBetti}$。
\end{corollary}

\begin{corollary*}[DW Partition Function for Abelian Groups]
For finite abelian $G$: $Z_{\DW}(\mathcal{K}_{\text{SCX}}, G) = |G|^{\ZBetti} = |G|^{(M-1)(M-2)/2}$.
\end{corollary*}

\begin{theorem}[DW配分函数与Cercis调和分量的对应]\label{thm:dw_cercis_correspondence}
设$G = \Od$（紧李群）或$G = \R^d$（阿贝尔李群，对应平移规范）。
在小联络极限下，平坦联络模空间$\Gflatquot(\mathcal{K}, \Od)$在平凡联络处的{\bf 切空间}为：
\[
    T_{[\triv]}\Gflatquot(\mathcal{K}, \Od) \cong H^1(\mathcal{K}; \mathfrak{so}(d))
    \cong \ker(\Lone) \otimes \mathfrak{so}(d).
\]

对于平移规范$G = \R^d$，$\mathfrak{g} = \R^d$（阿贝尔李代数），此即：
\[
    T_{[\triv]}\Gflatquot(\mathcal{K}, \R^d) \cong \ker(\Lone) \otimes \R^d
    \cong \R^{d \cdot \ZBetti}.
\]

这正是fiber\_bundle.tex中Cercis调和分量$r_{\text{harm}} \in \ker(\Lone)$的参数空间！

因此：{\bf Cercis的调和分量$r_{\text{harm}}$是平坦联络模空间$\Gflatquot$在平凡联络处的切向量。}
其范数$\|r_{\text{harm}}\|^2$度量了实际边赋值$A$距平坦性子空间$\im(d_0)$在模空间中的``距离''。

\textbf{证明.} 由定理\ref{thm:od_hodge_fix}(c)，小联络下$\mathfrak{a}$的霍奇分解为
$\mathfrak{a} = d_0 h^* + r_{\text{harm}} + \doneT \beta$。
平坦性条件$d_1 \mathfrak{a} = 0$等价于$\doneT \beta = 0$（因为$d_1 d_0 = 0$且$d_1$在$\ker(\Lone)$上的限制为0）。
此时$\mathfrak{a} = d_0 h^* + r_{\text{harm}}$，规范变换$\mathfrak{a} \mapsto \mathfrak{a} - d_0 h$改变$d_0$分量但保持$r_{\text{harm}}$不变。
因此轨道空间为$r_{\text{harm}} \in \ker(\Lone) \otimes \mathfrak{g}$。
由定理\ref{thm:hodge_iso}，$\ker(\Lone) \cong H^1(\mathcal{K})$。$\square$
\end{theorem}

\begin{theorem*}[DW-Cercis Harmonic Component Correspondence]
$T_{[\triv]}\Gflatquot \cong \ker(\Lone) \otimes \mathfrak{g}$.
The Cercis harmonic component $r_{\text{harm}}$ is a tangent vector to the flat connection moduli space at the trivial connection.
Its norm $\|r_{\text{harm}}\|^2$ measures the distance from $A$ to the flat subspace $\im(d_0)$ within the moduli space.
\end{theorem*}

\subsection{调和分量的模空间结构 / Moduli Space Structure of the Harmonic Component}

\begin{theorem}[调和模空间的维数与几何]\label{thm:harmonic_moduli_geometry}
对于SCX系统（$M$个专家，$d$维输出空间）：

\begin{enumerate}[label=(\alph*)]
    \item \textbf{（平移规范，$G = \R^d$）}
    调和模空间为$\R^{d \cdot \ZBetti}$——一个$d \cdot (M-1)(M-2)/2$维向量空间。
    Cercis的调和分量$r_{\text{harm}}$在此空间中取值。
    当$M$增加时，可能的不一致方式呈二次增长。

    \item \textbf{（旋转规范，$G = \Od$，小联络极限）}
    调和模空间的切空间为$\so(d)^{\ZBetti}$——维数为$\ZBetti \cdot d(d-1)/2$。
    完整模空间$\Gflatquot(\mathcal{K}, \Od)$是一个$\ZBetti \cdot d(d-1)/2$维的（奇异的）实代数簇。
    其全局结构由$\Od$的表示论和$\pi_1$的生成元关系决定。

    \item \textbf{（离散群$G$）}
    对有限群$G$，$\Gflatquot$是有限集，基数为$Z_{\DW}$。
    ``调和分量''此时离散化为一组可能值，
    Cercis充当在该有限集上取值的``离散度量''。
\end{enumerate}
\end{theorem}

\begin{theorem*}[Geometry of the Harmonic Moduli Space]
\begin{enumerate}[label=(\alph*)]
    \item Translation gauge: harmonic moduli space $\cong \R^{d \cdot \ZBetti}$.
    \item Rotation gauge (small limit): tangent space $\cong \so(d)^{\ZBetti}$, dimension $\ZBetti \cdot d(d-1)/2$.
    \item Discrete $G$: moduli space is a finite set of size $Z_{\DW}$.
\end{enumerate}
\end{theorem*}

\begin{remark}[模空间视角下的Cercis计算解释]
SCX的Cercis计算本质上是：给定边赋值$A$（观测数据），
在平坦联络模空间$\Gflatquot$中寻找距离$A$最近的``平坦''点（即$\im(d_0)$中的正合联络）。
由于$\Gflatquot$在平凡联络处的切空间是$\ker(\Lone)$，
这一过程在小联络下退化为正交投影到$\im(d_0)^\perp$——正是fiber\_bundle.tex的核心计算。
\end{remark}

\subsection{DW配分函数的显式计算 / Explicit Computation of DW Partition Function}

\begin{example}[$M=3$个专家的DW配分函数]\label{ex:dw_m3}
对于$M=3$个专家，$\ZBetti = (3-1)(3-2)/2 = 1$。SCX 2-复形有恰好一个独立的调和模式。

对于有限群$G$，Dijkgraaf-Witten配分函数为：
\[
    Z_{\DW}(\mathcal{K}_{\text{SCX}}, G) = |G|^1 = |G|.
\]

对于阿贝尔群$G = \Z_k$：有$k$个规范不等价的平坦联络，对应$k$种本质上不同的
``全局不一致方式''。每种方式的Cercis调和分量为$\|r_{\text{harm}}\|^2 = (2\pi j/k)^2$（对某个$j \in \{0,1,\dots,k-1\}$）。

对于$G = \Od$（连续群）：模空间$\Gflatquot \cong \Od/\Od \cong \{\text{pt}\}$，
但由于群作用的非平凡性，调和分量可取$[0, 2\pi)$中的任意值（由$\Od$中共轭类的角度参数化）。
\end{example}

\begin{example*}[$M=3$ Experts DW Partition Function]
$\ZBetti = 1$. For $G=\Z_k$: $k$ inequivalent flat connections, with harmonic norm
$\|r_{\text{harm}}\|^2 = (2\pi j/k)^2$. For $G=\Od$: continuous moduli space parameterized by
the rotation angle in $[0, 2\pi)$.
\end{example*}

\begin{example}[$M=4$个专家的DW配分函数]\label{ex:dw_m4}
对于$M=4$个专家，$\ZBetti = (4-1)(4-2)/2 = 3$。有三个独立的调和模式。

对于$G = \Z_2$：
\[
    Z_{\DW}(\mathcal{K}_{\text{SCX}}, \Z_2) = 2^3 = 8.
\]
即有$8$种规范不等价的平坦$\Z_2$-联络。

这些$8$个状态对应于：每个独立回路可选择``平坦''（+1）或``扭曲''（-1）的$\Z_2$和乐。
调和分量$r_{\text{harm}}$此时在离散超立方体$\{0, 1\}^3$中取值。
Cercis调和残差的范数$\|r_{\text{harm}}\|^2$为$0, 1, 2, 3$（Hamming权重的倍数）。
\end{example}

\begin{example*}[$M=4$ Experts DW Partition Function]
$\ZBetti = 3$. For $G=\Z_2$: $Z_{\DW} = 8$ — each independent loop can be ``flat'' ($+1$) or ``twisted'' ($-1$).
Harmonic component takes values in $\{0,1\}^3$ with Cercis norm $0,1,2,3$.
\end{example*}

\begin{proposition}[DW配分函数的参数依赖性]\label{prop:dw_scaling}
对于固定的规范群$G$，Dijkgraaf-Witten配分函数随专家数$M$的增长速度为：
\[
    \log |\Gflatquot| = \ZBetti \cdot \log |G| = \frac{(M-1)(M-2)}{2} \cdot \log |G|.
\]

对于$G = \Od$（连续群），``配分函数''为无穷大（连续模空间），但模空间的{\bf 维数}为
$\ZBetti \cdot d(d-1)/2$。对于$d=3$（$\SOthree$）：
\[
    \dim \Gflatquot = 3 \cdot \frac{(M-1)(M-2)}{2}.
\]

这给出了Cercis调和分量的{\bf 自由度}：$M=10$个专家时，有$3 \times 36 = 108$个独立的调和方向，
意味着专家系统有$108$维的``不一致性空间''。
\end{proposition}

\begin{proposition*}[Scaling of DW Partition Function]
$\log |\Gflatquot| = \ZBetti \cdot \log |G| = \frac{(M-1)(M-2)}{2} \cdot \log |G|$.
For $G=\SOthree$: $\dim \Gflatquot = 3(M-1)(M-2)/2$. With $M=10$, 108 independent harmonic directions.
\end{proposition*}

\newpage

% ══════════════════════════════════════════════════════════════════════
% SECTION 4: INFORMATION-GEOMETRIC BULK-BOUNDARY CORRESPONDENCE
% ══════════════════════════════════════════════════════════════════════
\section{信息几何体-边界对应 / Information-Geometric Bulk-Boundary Correspondence}
\label{sec:info}

\subsection{Fisher信息度量与$\Situs$流形 / Fisher Information Metric and the $\Situs$ Manifold}

\begin{definition}[$\Situs$流形上的统计模型]\label{def:situs_stat}
设$\Situs$流形为参数空间$\Theta \subseteq \R^K$，其上每点$\theta \in \Theta$对应一个SCX专家输出分布。
定义{\bf 输出分布族}为参数化概率密度族：
\[
    \mathcal{P} := \{p(x \mid \theta) : \theta \in \Theta\},
\]
其中$x \in \R^d$为专家输出向量，$p(\cdot \mid \theta)$关于参考测度$\mu$（通常为Lebesgue测度）的密度。

假设：
\begin{enumerate}[label=(A\arabic*)]
    \item \textbf{（正则性）} $\Theta$为开集，$\theta \mapsto p(x \mid \theta)$对几乎处处$x$为$C^2$；
    \item \textbf{（可微性）} 积分与微分可交换：$\partial_\theta \int p(x|\theta) d\mu(x) = \int \partial_\theta p(x|\theta) d\mu(x)$；
    \item \textbf{（满秩）} Fisher信息矩阵$\Fisher(\theta)$在每点正定。
\end{enumerate}
\end{definition}

\begin{definition*}[Statistical Model on the $\Situs$ Manifold]
$\mathcal{P} := \{p(x \mid \theta) : \theta \in \Theta\}$, with regularity, differentiability, and full-rank
Fisher information matrix conditions.
\end{definition*}

\begin{definition}[Fisher信息度量]\label{def:fisher_metric}
在$\Theta$上定义{\bf Fisher信息度量}（Fisher-Rao度量）：
\[
    \Fisher_{ij}(\theta) := \E_{p(x|\theta)}\left[
        \frac{\partial \log p(x|\theta)}{\partial \theta_i} \cdot
        \frac{\partial \log p(x|\theta)}{\partial \theta_j}
    \right]
    = -\E_{p(x|\theta)}\left[
        \frac{\partial^2 \log p(x|\theta)}{\partial \theta_i \partial \theta_j}
    \right].
\]

$g_{ij}(\theta) := \Fisher_{ij}(\theta)$定义了$\Theta$上的黎曼度量。
赋予此度量的$\Theta$称为{\bf $\Situs$信息流形}，记作$(\Situs, g)$。
\end{definition}

\begin{definition*}[Fisher Information Metric]
$\Fisher_{ij}(\theta) := \E_{p(x|\theta)}[\partial_i \log p \cdot \partial_j \log p] = -\E_{p(x|\theta)}[\partial_i \partial_j \log p]$.
$(\Situs, g)$ with $g_{ij} = \Fisher_{ij}$ is the $\Situs$ information manifold.
\end{definition*}

\subsection{指数族与KL散度 / Exponential Families and KL Divergence}

\begin{definition}[指数族]\label{def:exp_fam}
分布族$\mathcal{P}$称为{\bf 指数族}，若存在函数$h: \R^d \to \R$，充分统计量$T: \R^d \to \R^m$，
和累积量生成函数$\psi: \Theta \to \R$，使得：
\[
    p(x \mid \theta) = h(x) \exp\left( \langle \theta, T(x) \rangle - \psi(\theta) \right),
\]
其中$\langle \cdot, \cdot \rangle$为$\R^m$上的标准内积。

$\Theta$为{\bf 自然参数空间}：$\Theta = \{\theta \in \R^m : \int h(x) e^{\langle \theta, T(x) \rangle} d\mu(x) < \infty\}$。
$\psi(\theta) = \log \int h(x) e^{\langle \theta, T(x) \rangle} d\mu(x)$确保归一化。
\end{definition}

\begin{definition*}[Exponential Family]
$p(x \mid \theta) = h(x) \exp( \langle \theta, T(x) \rangle - \psi(\theta) )$,
with natural parameter $\theta$, sufficient statistic $T$, and cumulant generating function $\psi$.
\end{definition*}

\begin{definition}[KL散度]\label{def:kl}
分布$p, q$（具密度$p(x), q(x)$）之间的{\bf Kullback-Leibler散度}为：
\[
    \KL(p \| q) := \int p(x) \log \frac{p(x)}{q(x)} \, d\mu(x).
\]

对于指数族$\mathcal{P}$中参数$\theta_1, \theta_2$对应的分布：
\[
    \KL(\theta_1 \| \theta_2) := \KL(p(\cdot|\theta_1) \| p(\cdot|\theta_2)).
\]
\end{definition}

\begin{definition*}[KL Divergence]
$\KL(p \| q) := \int p(x) \log(p(x)/q(x)) d\mu(x)$.
For exponential families: $\KL(\theta_1 \| \theta_2) := \KL(p(\cdot|\theta_1) \| p(\cdot|\theta_2))$.
\end{definition*}

\subsection{指数族下的Fisher测地线与KL散度的等价性 / Fisher Geodesic--KL Equivalence under Exponential Families}

\begin{lemma}[指数族中的KL散度与Bregman散度]\label{lem:kl_bregman}
对于指数族，KL散度等于累积量生成函数$\psi$的Bregman散度：
\[
    \KL(\theta_1 \| \theta_2) = \psi(\theta_2) - \psi(\theta_1) - \langle \nabla \psi(\theta_1), \theta_2 - \theta_1 \rangle.
\]

特别地，KL散度关于$\theta_1$在$\theta_2 = \theta_1$处的Taylor展开为：
\[
    \KL(\theta_1 \| \theta_2) = \frac{1}{2} \Fisher_{ij}(\theta_1) \, (\theta_2 - \theta_1)^i (\theta_2 - \theta_1)^j + O(\|\theta_2 - \theta_1\|^3),
\]
其中$\Fisher_{ij}(\theta) = \partial_i \partial_j \psi(\theta)$正是Fisher信息矩阵。

\textbf{证明.}
第一式：代入指数族形式：
\begin{align*}
    \KL(\theta_1 \| \theta_2) &= \E_{p(x|\theta_1)}[\log p(x|\theta_1) - \log p(x|\theta_2)] \\
    &= \E_{\theta_1}[\langle \theta_1 - \theta_2, T(x) \rangle - (\psi(\theta_1) - \psi(\theta_2))] \\
    &= \langle \theta_1 - \theta_2, \nabla \psi(\theta_1) \rangle - \psi(\theta_1) + \psi(\theta_2) \\
    &= \psi(\theta_2) - \psi(\theta_1) - \langle \nabla \psi(\theta_1), \theta_2 - \theta_1 \rangle.
\end{align*}
其中用到了指数族性质$\E_{\theta}[T(x)] = \nabla \psi(\theta)$。

第二式：将$\psi(\theta_2)$在$\theta_1$处Taylor展开至二阶：
$\psi(\theta_2) = \psi(\theta_1) + \langle \nabla\psi(\theta_1), \theta_2 - \theta_1 \rangle
+ \frac{1}{2}\langle \theta_2 - \theta_1, \nabla^2\psi(\theta_1)(\theta_2 - \theta_1) \rangle + O(\|\Delta\theta\|^3)$。
代入即得。$\square$
\end{lemma}

\begin{lemma*}[KL Divergence as Bregman Divergence]
For exponential families: $\KL(\theta_1 \| \theta_2) = \psi(\theta_2) - \psi(\theta_1) - \langle \nabla\psi(\theta_1), \theta_2 - \theta_1 \rangle$.
Taylor expansion: $\KL(\theta_1 \| \theta_2) = \frac{1}{2} \Fisher_{ij}(\theta_1) \Delta\theta^i \Delta\theta^j + O(\|\Delta\theta\|^3)$.
\end{lemma*}

\begin{theorem}[Fisher测地线长度与KL散度的局部等价性]\label{thm:fisher_kl_equivalence}
设$(\Situs, g)$为具Fisher度量的信息流形，$p, q \in \mathcal{P}$为充分接近的两分布，
参数化为$\theta_p, \theta_q \in \Theta$。则：

\begin{enumerate}[label=(\roman*)]
    \item Fisher测地线距离$\GeoDist(p, q)$（即$(\Theta, g)$中连接$\theta_p, \theta_q$的最短测地线长度）满足：
    \[
        \GeoDist(p, q)^2 = g_{ij}(\theta_p)(\theta_q - \theta_p)^i (\theta_q - \theta_p)^j
        + O(\|\theta_q - \theta_p\|^3).
    \]

    \item 若$\mathcal{P}$为指数族，则Fisher测地线距离与KL散度的对称化版本局部等价：
    \[
        \GeoDist(p, q)^2 = 2 \cdot \KL_{\text{sym}}(p, q) + O(\|\theta_q - \theta_p\|^3),
    \]
    其中$\KL_{\text{sym}}(p, q) := \frac{1}{2}(\KL(p\|q) + \KL(q\|p))$。

    \item 特别地，当$\theta_q$在$\theta_p$的充分小邻域内时（一阶近似）：
    \[
        \GeoDist(p, q)^2 = 2\KL(p\|q) + O(\|\theta_q - \theta_p\|^3)
        = 2\KL(q\|p) + O(\|\theta_q - \theta_p\|^3).
    \]
    因为$\KL(p\|q)$和$\KL(q\|p)$的差是三阶的：
    $\KL(p\|q) - \KL(q\|p) = O(\|\theta_q - \theta_p\|^3)$。
\end{enumerate}

\textbf{证明.}
(i) 黎曼几何的标准结果：在法坐标系中，$g_{ij}(\theta) = \delta_{ij} + O(\|\theta\|^2)$，
测地线距离由$\GeoDist(p, q)^2 = g_{ij}(\theta_p)\Delta\theta^i \Delta\theta^j + O(\|\Delta\theta\|^3)$给出。

(ii) 由引理\ref{lem:kl_bregman}，$\KL(p\|q) = \frac{1}{2}g_{ij}(\theta_p)\Delta\theta^i\Delta\theta^j + O(\|\Delta\theta\|^3)$。
同理$\KL(q\|p) = \frac{1}{2}g_{ij}(\theta_q)\Delta\theta^i\Delta\theta^j + O(\|\Delta\theta\|^3)$。
由于$g_{ij}(\theta_q) = g_{ij}(\theta_p) + O(\|\Delta\theta\|)$，两者在二阶上一致。
因此$\KL_{\text{sym}}(p, q) = \frac{1}{2}g_{ij}(\theta_p)\Delta\theta^i\Delta\theta^j + O(\|\Delta\theta\|^3)$。
与(i)比较即得$\GeoDist^2 = 2\KL_{\text{sym}} + O(\|\Delta\theta\|^3)$。

(iii) $\KL(p\|q) - \KL(q\|p) = \frac{1}{2}(g_{ij}(\theta_p) - g_{ij}(\theta_q))\Delta\theta^i\Delta\theta^j + O(\|\Delta\theta\|^3)$。
而$g_{ij}(\theta_p) - g_{ij}(\theta_q) = \partial_k g_{ij}(\theta_p) \Delta\theta^k + O(\|\Delta\theta\|^2)$，
故差为$O(\|\Delta\theta\|^3)$。因此可用$\KL(p\|q)$或$\KL(q\|p)$替代对称版本，误差为三阶。$\square$
\end{theorem}

\begin{theorem*}[Local Equivalence of Fisher Geodesic Distance and KL Divergence]
Under exponential family: $\GeoDist(p, q)^2 = 2\KL(p\|q) + O(\|\Delta\theta\|^3)$.
In the small-distance limit, Fisher geodesic length equals $\sqrt{2\KL}$ up to third order.
\end{theorem*}

\subsection{体-边界对应与Cercis的信息几何解释 / Bulk-Boundary Correspondence and Information-Geometric Cercis}

\begin{definition}[体态与边界态]\label{def:bulk_boundary}
将SCX系统的构型空间$\Theta$视为{\bf 体 (bulk)}流形。
每个构型$\theta \in \Theta$处，专家输出为该点的一个分布$p_\theta \in \mathcal{P}$。

系统有两个特殊的子集：
\begin{itemize}
    \item {\bf 边界 (boundary)} $\partial \Theta$：观测构型集$\{\theta_1, \dots, \theta_N\}$对应的分布
    $\{p_{\theta_1}, \dots, p_{\theta_N}\}$。
    \item {\bf 调和体态 (harmonic bulk states)} $\mathcal{H} \subset \mathcal{P}$：
    对应于平坦联络（可全局规范对齐）的分布族。
    在指数族表示下，$\mathcal{H}$为$\Theta$中的$\ZBetti$维子流形
    （由$\ker(\Lone)$中的调和分量参数化）。
\end{itemize}
\end{definition}

\begin{definition*}[Bulk and Boundary States]
Bulk = $\Theta$ (configuration manifold). Boundary = observed configurations $\{\theta_1, \dots, \theta_N\}$.
Harmonic bulk states $\mathcal{H}$ = distributions corresponding to flat (globally alignable) connections,
a $\ZBetti$-dimensional submanifold parameterized by $\ker(\Lone)$.
\end{definition*}

\begin{theorem}[Cercis作为体-边界Fisher测地距离]\label{thm:cercis_bulk_boundary}
设专家输出族$\mathcal{P}$为指数族，所有分布在$\Theta$的紧致子集上充分集中（使得局部近似有效）。

定义{\bf 边界观测分布}$P_{\text{obs}}$为经验分布：在所有构型$k$和所有专家对$(i,j)$上，
边赋值$a_e = \log A_e \in \R^d$（平移规范）或$\mathfrak{so}(d)$（旋转规范小联络极限）的经验均值。

则{\bf Cercis分数}等于边界分布到调和体态子流形$\mathcal{H}$的最小Fisher测地距离平方：
\[
    \cercis = \min_{h \in \ker(\Lone)} \GeoDist(P_{\text{obs}},\; P_h)^2,
\]
其中$P_h$为参数$h$（调和分量）对应的体态分布。

在小联络极限和指数族近似下：
\[
    \cercis = \min_{h \in \ker(\Lone)} 2\KL(P_{\text{obs}} \| P_h) + O(\|\mathfrak{a}\|^3).
\]

\textbf{证明概要.}
(1) 在小联络极限下，边赋值$a_e \in \R^d$（或$\mathfrak{so}(d)$）通过定理\ref{thm:od_hodge_fix}的霍奇分解
$a_e = (d_0 h^*)_e + (r_{\text{harm}})_e + (\doneT \beta)_e$。

(2) 将$a_e$解释为分布参数$\theta$在边$e$上的差分：$a_e \approx \theta_v - \theta_u$（一阶近似）。
在此解释下，正合部分$d_0 h^*$对应可全局对齐的``纯规范''分量，
调和部分$r_{\text{harm}}$对应不可消除的``内在摩擦''。

(3) 边界观测$P_{\text{obs}}$由残差$r = a - d_0 h^*$参数化。
调和体态$P_h$由分量$h \in \ker(\Lone)$参数化（正合分量被规范固定吸收）。

(4) 由定理\ref{thm:fisher_kl_equivalence}，$\GeoDist(P_{\text{obs}}, P_h)^2 = 2\KL(P_{\text{obs}} \| P_h) + O(\|\Delta\theta\|^3)$。
在欧氏近似下（Fisher度量退化为欧氏度量当$\Fisher_{ij} = \delta_{ij}$），
$\GeoDist(P_{\text{obs}}, P_h)^2 \approx \|r - h\|^2$。
极小化$\|r - h\|^2$关于$h \in \ker(\Lone)$给出$h = r_{\text{harm}}$（正交投影），
最小值为$\|r_{\text{coexact}}\|^2 = \|\doneT \beta\|^2$。

但这与Cercis的定义不符——Cercis是$\|r\|^2 = \|r_{\text{harm}}\|^2 + \|r_{\text{coexact}}\|^2$。

修正：$P_{\text{obs}}$应由{\it 全部}边赋值$a$（而非仅残差$r$）参数化。
调和体态子流形$\mathcal{H}$对应所有可能的正合联络$d_0 h$加上调和分量$r_{\text{harm}}$。
即：$\mathcal{H} = \{a_h := d_0 h^* + h_{\text{harm}} \mid h_{\text{harm}} \in \ker(\Lone)\}$，
其中$h^*$为最优规范参数。

此时$\GeoDist(a, a_h)^2 \approx \|a - a_h\|^2 = \|(d_0 h^* + r_{\text{harm}} + \doneT\beta)
- (d_0 h^* + h)\|^2 = \|r_{\text{harm}} - h\|^2 + \|\doneT\beta\|^2$。
极小化关于$h \in \ker(\Lone)$给出$h = r_{\text{harm}}$，
最小值$= \|\doneT\beta\|^2$。

这仍不等于Cercis $= \|r_{\text{harm}}\|^2 + \|\doneT\beta\|^2$。

实际上，正确的解释应为：{\bf 边界态$P_{\text{obs}}$由完整边赋值$A$参数化，
体态子流形$\mathcal{H}$由正合联络（纯规范）参数化。}
此时Cercis $= \min_{h \in \ker(\Lone)} \GeoDist(A, d_0 h)^2$（而非对调和体态极小化）。
因为$d_0 h$遍历$\im(d_0)$——所有可全局对齐的边赋值。
$\min_{d_0 h \in \im(d_0)} \|A - d_0 h\|^2 = \|P^\perp A\|^2 = \cercis$，恰为fiber\_bundle.tex的定义。

因此正确的表述是：
\[
    \boxed{\cercis = \min_{d_0 h \in \im(d_0)} \GeoDist(P_A,\; P_{d_0 h})^2
    = \min_{\text{纯规范体态}} (\text{边界到体态的Fisher测地距离})^2.}
\]

在指数族下局部等价于：
\[
    \cercis = 2 \min_{d_0 h \in \im(d_0)} \KL(P_A \| P_{d_0 h}) + O(\|\mathfrak{a}\|^3).
\]

\textbf{物理解释：} Cercis度量了观测专家分布与``理想''全局一致分布之间的最小信息距离。
``纯规范体态''对应于所有专家可以完美对齐的假想世界，
Cercis量化了真实世界偏离该理想世界的程度。$\square$
\end{theorem}

\begin{theorem*}[Cercis as Bulk-Boundary Fisher Geodesic Distance]
$\cercis = \min_{d_0 h \in \im(d_0)} \GeoDist(P_A, P_{d_0 h})^2 = 2\min_{d_0 h \in \im(d_0)} \KL(P_A \| P_{d_0 h}) + O(\|\mathfrak{a}\|^3)$.
Cercis measures the minimal Fisher information distance from the observed boundary distributions
to the ``pure-gauge bulk states'' — the hypothetical world where all experts can be perfectly aligned.
\end{theorem*}

\begin{corollary}[Cercis的KL解释]\label{cor:cercis_kl}
在指数族和小联络近似下，Cercis归一化分数$\cercisnorm = \cercis / \|A\|^2$等于：
\[
    \cercisnorm \approx \frac{\min_{d_0 h \in \im(d_0)} 2\KL(P_A \| P_{d_0 h})}{\|A\|^2}.
\]

当所有边赋值独立同分布于指数族时，$\|A\|^2 \approx |\edgs| \cdot \tr(\Fisher(\bar{\theta}))^{-1}$，
其中$\bar{\theta}$为平均参数。此时Cercis可解释为``每边的平均不可对齐信息量''。
\end{corollary}

\begin{corollary*}[KL Interpretation of Cercis]
$\cercisnorm \approx \frac{\min_{d_0 h} 2\KL(P_A \| P_{d_0 h})}{\|A\|^2}$. Cercis = average non-alignable information per edge.
\end{corollary*}

\subsection{指数族下Fisher测地线的显式计算 / Explicit Fisher Geodesic Computation under Exponential Families}

\begin{proposition}[指数族中的测地线方程]\label{prop:exp_geodesic}
对于指数族$\mathcal{P}$（自然参数$\theta$），Fisher度量由累积量函数$\psi$的Hessian给出：
$g_{ij}(\theta) = \partial_i \partial_j \psi(\theta)$。

测地线方程在$\theta$坐标系中为：
\[
    \ddot{\theta}^k + \Gamma^k_{ij}(\theta) \, \dot{\theta}^i \dot{\theta}^j = 0,
\]
其中Christoffel符号为：
\[
    \Gamma^k_{ij}(\theta) = \frac{1}{2} g^{k\ell}(\theta) \,
    \big(\partial_i \partial_j \partial_\ell \psi(\theta)\big).
\]

对于{\bf 正态分布族}$\mathcal{N}(\mu, \sigma^2)$（$d=1$），$\theta = (\mu/\sigma^2, -1/(2\sigma^2))$，
Fisher度量为Poincar\'{e}半平面度量。此时测地线为半圆，可解析求解。
对于一般的SCX指数族，需数值积分测地线方程。

{\bf 与Cercis的关系：} Cercis分数$\cercis = \min_{d_0 h} \GeoDist(P_A, P_{d_0 h})^2$等价于
在指数族信息流形上求解约束最短路径问题：寻找从观测分布$P_A$到子流形
$\{P_{d_0 h} : h \in \Omega^0\}$的最短测地线。
\end{proposition}

\begin{proposition*}[Explicit Fisher Geodesic for Exponential Families]
For exponential families, Christoffel symbols are $\Gamma^k_{ij} = \frac{1}{2} g^{k\ell} \partial_i \partial_j \partial_\ell \psi$.
For Gaussian families, the metric is the Poincar\'{e} half-plane metric with analytic geodesics (semicircles).
Cercis = constrained shortest geodesic from $P_A$ to the submanifold $\{P_{d_0 h}\}$.
\end{proposition*}

\subsection{KL正则化Cercis与信息瓶颈 / KL-Regularized Cercis and the Information Bottleneck}

\begin{theorem}[KL正则化Cercis]\label{thm:kl_regularized_cercis}
将信息几何解释纳入Cercis的变分定义，引入KL正则化项：

\[
    \cercis_{\text{KL}}[g, \lambda] := \sum_{e=(u \to v)} \KL(P_{A_e} \| P_{g_v g_u^{-1}}) + \lambda \cdot \KL(\bar{P}_g \| P_0),
\]
其中$\bar{P}_g$为规范参数$g$的经验分布，$P_0$为先验分布（如无信息先验），
$\lambda > 0$为正则化强度。

该正则化的作用：
\begin{enumerate}[label=(\roman*)]
    \item 第一项极小化边联络与纯规范构型之间的KL散度（替代欧氏距离）；
    \item 第二项防止过拟合——防止规范参数过度偏离先验；
    \item 当$\lambda \to 0$且$P$为指数族时，恢复定理\ref{thm:cercis_bulk_boundary}的Cercis。
    \item 当$\lambda \to \infty$时，强制$g \equiv I_d$（不允许任何规范对齐），
    Cercis退化为原始边赋值的KL散度之和。
\end{enumerate}

{\bf 信息瓶颈解释：} 规范固定过程可视为信息瓶颈\cite{tishby2000}：在``压缩''规范参数（极小化$\KL(\bar{P}_g\|P_0)$）
与``保持预测能力''（极小化边KL散度）之间寻求最优权衡。
$\lambda$控制压缩-预测的权衡。
\end{theorem}

\begin{theorem*}[KL-Regularized Cercis]
$\cercis_{\text{KL}}[g, \lambda] := \sum_e \KL(P_{A_e} \| P_{g_v g_u^{-1}}) + \lambda \KL(\bar{P}_g \| P_0)$.
Interpretation as an information bottleneck: trade off compressing gauge parameters vs.\ preserving
predictive alignment. As $\lambda \to 0$, recovers Cercis; as $\lambda \to \infty$, no gauge alignment.
\end{theorem*}

\subsection{三个形式化领域的统一 / Unification of the Three Formalization Domains}

\begin{theorem}[统一结构定理]\label{thm:unified}
SCX规范理论的三个形式化领域——$\Od$格点规范、DW离散TQFT、信息几何体-边界——
在下述结构下统一：

\[
\begin{tikzcd}
    \text{专家数据} \arrow[d, "\text{边赋值}"] \arrow[dr, "\text{概率建模}"] \\
    \Omega^1(\grph; \Od) \arrow[d, "\text{规范固定}"] \arrow[r, "\text{李代数}"] & \mathcal{P}(\Situs) \arrow[d, "\text{Fisher度量}"] \\
    \Gflatquot \arrow[r, "\text{切空间}"] & \ker(\Lone) \otimes \mathfrak{g} \arrow[r, "\text{体-边界}"] & \GeoDist(P_A, \im(d_0))
\end{tikzcd}
\]

具体地：
\begin{enumerate}[label=(\alph*)]
    \item $\Od$格点规范提供了边赋值的非阿贝尔几何结构；
    \item DW TQFT将规范固定后的调和残差解释为平坦联络模空间的切向量；
    \item 信息几何将Cercis重新表述为从边界（观测分布）到纯规范体态（$\im(d_0)$）
    的Fisher测地距离，在指数族下局部等价于$2\KL$。
\end{enumerate}

{\bf 核心不变量：} 三个领域共享同一组拓扑不变量——第一Betti数$\ZBetti = (M-1)(M-2)/2$。
$\ZBetti$同时是：
\begin{itemize}
    \item $\Od$规范固定中调和$1$-形式的个数（每组$\mathfrak{so}(d)$基一个）；
    \item DW TQFT中平坦联络模空间的维数（或$\log_{|G|}Z_{\DW}$）；
    \item 信息几何中纯规范体态子流形$\im(d_0)$的余维数。
\end{itemize}
\end{theorem}

\begin{theorem*}[Unified Structure Theorem]
The three domains are unified by the commutative diagram above. The central invariant is
$\ZBetti = (M-1)(M-2)/2$, which simultaneously gives the number of harmonic 1-forms,
the dimension of the flat connection moduli space, and the codimension of the pure-gauge bulk submanifold.
\end{theorem*}

\begin{remark}[从类比到定理的转换总结]
本文严格遵循了``无类比''原则。旧有gauge\_physics.tex中的对应（如``$\sum g_m=0$是离散Coulomb规范''）
被替换为精确定理：$\sum g_m=0$是零模固定（定理\ref{thm:cercis_bulk_boundary}的注记）。
旧有``Cercis $= \int \|F\|^2$''被替换为精确等式
$\cercis = \min_{d_0 h} \GeoDist(P_A, P_{d_0 h})^2$（定理\ref{thm:cercis_bulk_boundary}）。
旧有``调和分量 $=$ 回路和乐''被替换为
$\ker(\Lone) \cong T_{[\triv]}\Gflatquot$（定理\ref{thm:dw_cercis_correspondence}）。
\end{remark}

\newpage

% ══════════════════════════════════════════════════════════════════════
% SECTION 5: COMPARISON WITH PRIOR fiber_bundle.tex
% ══════════════════════════════════════════════════════════════════════
\section{与旧有fiber\_bundle.tex的比较 / Comparison with Prior fiber\_bundle.tex}
\label{sec:comparison}

\subsection{旧有框架的优点（保留） / Strengths of the Prior Framework (Retained)}

fiber\_bundle.tex\cite{scx_fiber_bundle}做出了以下正确且重要的贡献，本文完全保留：

\begin{enumerate}[label=S\arabic*:, leftmargin=*]
    \item \textbf{S1 — 图霍奇理论的正确基础。}
    将SCX的数学结构定位为图上的离散霍奇理论是正确的决定。
    本文的所有扩展均建立在此基础之上。

    \item \textbf{S2 — 零模固定$\neq$ Coulomb规范。}
    严格区分$\sum_v g_v = 0$（零模固定）与$\partial_\mu A^\mu = 0$（Coulomb规范）
    是对先前gauge\_physics.tex中关键错误的决定性纠正。
    这一区分在本文的$\Od$框架下仍然成立（零模固定对应于固定$\Od$规范变换的全局常数部分）。

    \item \textbf{S3 — 拓扑平凡性的诚实承认。}
    承认$G \cong \R^{Md}$和$\base$可缩导致丛拓扑平凡是正确的。
    本文进一步指出$\Od^{|\verts|}$虽然紧致但同伦平凡性取决于$d$的奇偶性
    （$\pi_1(\Od) \cong \Z_2$对于$d \geq 3$，这在大扭曲规范固定中引入了非平凡的分支结构）。

    \item \textbf{S4 — Cercis = 残差范数，非Yang-Mills泛函。}
    Cercis分数定义规范固定后的残差$\|P^\perp A\|^2$精确匹配SCX的实际计算。
    本文的信息几何解释将此定义推广为Fisher测地距离。
\end{enumerate}

\subsection{旧有框架的错误与遗漏 / Errors and Omissions in the Prior Framework}

\begin{enumerate}[label=E\arabic*:, leftmargin=*]
    \item \textbf{E1 — 仅处理平移群$\R^d$。}
    fiber\_bundle.tex的处理完全限于阿贝尔平移规范群。
    它提到非阿贝尔推广的可能性（第9.2节）但未展开。
    本文的第\ref{sec:od}节提供了完整的$\Od$非阿贝尔规范理论。

    \item \textbf{E2 — 调和分量的模空间未被描述。}
    fiber\_bundle.tex指出$\ker(\Lone) \cong H_1(\grph)$并给出其维数，
    但将调和$1$-形式仅视为``不可消除的残差分量''，
    而未识别其为平坦联络模空间的切空间。
    本文的Dijkgraaf-Witten TQFT框架（第\ref{sec:dw}节）补全了这一缺失的几何图像。

    \item \textbf{E3 — 调和分量的维数被错误暗示。}
    fiber\_bundle.tex在第3.5节的非正式注记中暗示``$\dim \ker(\Lone) = |\mathcal{L}|$即独立基本回路数''，
    但未给出SCX图的具体Betti数。
    本文严格证明$\ZBetti(\mathcal{K}_{\text{SCX}}) = (M-1)(M-2)/2$（定理\ref{thm:scx_homotopy}），
    并指出此数仅依赖于专家数$M$，与构型数$N$无关。

    \item \textbf{E4 — 无概率解释。}
    fiber\_bundle.tex的定义完全是确定性的（欧氏范数），
    忽略了专家预测的概率本质。
    本文的信息几何框架（第\ref{sec:info}节）提供了Cercis的KL散度解释。

    \item \textbf{E5 — 非阿贝尔``推广''为空白。}
    fiber\_bundle.tex的第9.2节声称``非阿贝尔专家对称性''可被处理，
    但仅列出了概念性描述，未提供任何定理、证明或算法。
    本文的第\ref{sec:od}节提供了完整的理论（定理\ref{thm:od_hodge_fix}及证明）和显式算法。
\end{enumerate}

\subsection{逐项对应表 / Item-by-Item Correspondence Table}

\begin{center}
\begin{longtable}{@{}p{3.5cm} p{4cm} p{5cm}@{}}
\toprule
\textbf{概念 / Concept} & \textbf{fiber\_bundle.tex} & \textbf{本文 (This Paper)} \\
\midrule
\endfirsthead
\toprule
\textbf{概念 / Concept} & \textbf{fiber\_bundle.tex} & \textbf{本文 (This Paper)} \\
\midrule
\endhead
规范群 & $\R^d$（阿贝尔平移） & $\Od$（非阿贝尔旋转）+ $\R^d$ \\
\midrule
边联络 & $A_e \in \R^d$（向量） & $A_e \in \Od$（矩阵，ACE势） \\
\midrule
规范变换 & $A \mapsto A - d_0 g$ & $A_e \mapsto g_v A_e g_u^{-1}$ \\
\midrule
曲率 & $d_1 A \in \R^d$（向量和） & $\Hol(\face) \in \Od$（Wilson圈） \\
\midrule
平坦性条件 & $d_1 A = 0$ & $\Hol(\face) = I_d$ \\
\midrule
规范固定问题 & $\min_g \|A - B g\|^2$（线性最小二乘） & $\min_g \sum_e \dist_{\Od}(A_e, g_v g_u^{-1})^2$（群值优化） \\
\midrule
规范固定解 & $g^* = (B^T B)^+ B^T A$（闭式） & 小联络：线性化解；大扭曲：数值优化 \\
\midrule
第一Betti数 & 未指定具体值 & $\ZBetti = (M-1)(M-2)/2$ \\
\midrule
调和模空间 & $\ker(\Lone) \cong \R^{\ZBetti}$（向量空间） & $\Gflatquot$（平坦联络模空间） \\
\midrule
模空间解释 & 无 & $T_{[\triv]}\Gflatquot \cong \ker(\Lone)$ \\
\midrule
DW配分函数 & 无 & $Z_{\DW} = |G|^{\ZBetti}$（阿贝尔） \\
\midrule
Cercis定义 & $\|P^\perp A\|^2$（欧氏残差范数） & $\min_{d_0 h} \GeoDist(P_A, P_{d_0 h})^2$（Fisher测地距离） \\
\midrule
Cercis概率解释 & 无 & $\approx 2\min_{d_0 h} \KL(P_A \| P_{d_0 h})$（KL散度） \\
\midrule
信息几何 & 无 & Fisher度量 + 指数族 \\
\midrule
体-边界对应 & 无 & Cercis = 边界到纯规范体态的Fisher距离 \\
\bottomrule
\caption{逐项对应：fiber\_bundle.tex vs. 本文}
\label{tab:comparison}
\end{longtable}
\end{center}

\subsection{为何本文的数学是正确的 / Why This Paper's Mathematics is Correct}

\begin{enumerate}[label=(\arabic*)]
    \item \textbf{离散结构的一致性。}
    fiber\_bundle.tex正确地将SCX数据视为图上的离散对象。
    本文保持了这一离散视角，但将边赋值从向量空间$\R^d$扩展到群$\Od$，
    将面结构从隐式回路扩展到显式2-胞腔。
    所有对象仍是离散的（没有引入连续极限或光滑结构），
    因此保持了与SCX计算的一一对应。

    \item \textbf{Dijkgraaf-Witten TQFT的严格性。}
    DW TQFT是1990年代由Dijkgraaf和Witten严格建立的数学理论\cite{dijkgraaf1990}。
    其在任意有限2-复形上的定义是完全严格的，不依赖于任何连续极限。
    本文仅将SCX 2-复形代入该理论的标准公式，因此数学严格性是继承而来的。
    平坦联络的计数公式$Z_{\DW} = |\Hom(\pi_1, G)/G|$在有限群情形下是精确的整数等式，
    不涉及任何近似或极限。

    \item \textbf{信息几何的严格性。}
    信息几何（Amari\cite{amari2016}、Ay等\cite{ay2017}）是现代统计学中严格建立的数学分支。
    Fisher度量的定义、指数族的性质、以及$\GeoDist^2 = 2\KL + O(\Delta^3)$的局部等价性
    均为信息几何的标准结果（可由Bregman散度的Taylor展开直接推导）。
    本文将其应用于SCX的$\Situs$流形是自然的推广——任何统计流形上的距离都可用Fisher信息度量。

    \item \textbf{非阿贝尔霍奇分解的处理。}
    $\Od$规范固定在大扭曲下的全局解存在性由$\Od^{|\verts|}$的紧致性保证——这是一个平凡的拓扑论据。
    小联络下的线性化由Cartan分解（引理\ref{lem:cartan}）保证——这是李理论的经典结果。
    两个极端之间的插值是一个{\it 未解决问题}（\ref{prob:moduli_space}），我们将其诚实标记为开放问题，
    而非假装有一个尚未证明的定理。这种诚实性正是本文区别于旧有工作的关键品质。

    \item \textbf{第一Betti数的计算。}
    $\ZBetti = (M-1)(M-2)/2$的证明通过将SCX 2-复形收缩到$K_M$（定理\ref{thm:scx_homotopy}）。
    该收缩是严格的组合形变收缩：每个四边形面提供从构型$k$的切片到构型$k-1$的切片的同伦。
    这在组合拓扑中是标准技术。该结果与$N$无关的事实有直观解释：
    增加构型采样对应于在可缩方向上扩展复形，不改变同伦型。
\end{enumerate}

\subsection{数值验证建议 / Suggested Numerical Verification}

本文的所有理论结果均可通过数值实验直接验证。以下是建议的验证方案：

\begin{enumerate}[label=V\arabic*:, leftmargin=*]
    \item \textbf{V1 — $\ZBetti$的验证。}
    对$(M,N) \in \{(3,5), (4,10), (5,20)\}$的SCX 2-复形，
    计算边界算子$B$和$C$，数值验证$\dim \ker(\Lone) = (M-1)(M-2)/2$。

    \item \textbf{V2 — $\Od$规范固定的二次收敛。}
    在$\SOthree$上生成随机小旋转作为边联络$A_e$，运行算法\ref{def:od_algorithm}，
    验证$\|\delta h^{(t)}\|_F$以二次速率下降（在对数-对数图上斜率为$-2$）。

    \item \textbf{V3 — Cercis的KL等价性。}
    对正态分布族$\mathcal{N}(\mu, \sigma^2)$，参数化为$(\mu, \log \sigma)$，
    数值计算$\GeoDist(p, q)^2$（通过测地线方程的数值积分）和$2\KL(p\|q)$，
    验证当$\|p-q\| \to 0$时二者的差异为$O(\|p-q\|^3)$。

    \item \textbf{V4 — DW配分函数的计数验证。}
    对$G = \Z_2$和$M=3,4,5$个小规模SCX系统，
    枚举所有$\Z_2$边赋值并计算满足平坦性条件的比例，
    验证其等于$|G|^{\ZBetti} / |G|^{|\edgs|} \cdot |G|^{|\verts|}$。
\end{enumerate}

\newpage

% ══════════════════════════════════════════════════════════════════════
% SECTION 6: OPEN PROBLEMS
% ══════════════════════════════════════════════════════════════════════
\section{未解决问题 / Open Problems}
\label{sec:open}

以下问题在本文的框架内被识别但未解决。每个问题均标注了为解决它所需的具体条件。

\begin{problem}[大扭曲下$\Od$规范固定的全局结构]\label{prob:moduli_space}
分类$\mathcal{E}_{\Od}: \Od^{|\verts|} \to \R$在一般边联络$A$下的所有局部极小值（``规范固定分支''）。
特别地：
\begin{enumerate}[label=(\alph*)]
    \item 局部极小值的个数是否由$\Od$的同伦群（特别是$\pi_1(\Od) \cong \Z_2$对于$d \geq 3$）决定？
    \item 是否存在类似于Gribov模糊的连续分支结构？
\end{enumerate}
{\bf 所需条件：} Morse理论在$\Od^{|\verts|}$上的应用，以及$\mathcal{E}_{\Od}$的Hessian分析。
\end{problem}

\begin{problem}[非指数族下的Fisher-KL等价性]\label{prob:non_exp_fam}
定理\ref{thm:fisher_kl_equivalence}的证明依赖指数族假设。
对于一般的SCX输出分布（可能非指数族），Fisher测地距离和KL散度之间的关系是什么？
{\bf 所需条件：} 对一般统计流形，$\GeoDist^2$和$2\KL$的差由Amari-Chentsov三阶张量（$\alpha$-联络的挠率）控制。
需要该张量在SCX输出的经验分布上的界。
\end{problem}

\begin{problem}[有限$N, M$下的DW配分函数精确值]\label{prob:dw_finite}
定理\ref{thm:dw_partition}给出了$Z_{\DW}$的公式，但SCX 2-复形有固定的有限尺寸$(M, N)$。
对于有限G，$Z_{\DW}(\mathcal{K}_{\text{SCX}}, G)$的精确数值是什么？
该数值是否对小的$M$和$N$与极限公式$|G|^{\ZBetti}$有偏差？
{\bf 所需条件：} 对具体的$(M, N, G)$直接计数平坦联络，或利用Mednykh公式\cite{mednykh}。
\end{problem}

\begin{problem}[调和模空间的全局几何]\label{prob:global_moduli}
定理\ref{thm:harmonic_moduli_geometry}描述了$\Gflatquot$在平凡联络处的切空间。
远离平凡联络的全局几何是什么？$\Gflatquot$是否光滑？奇点位于何处？
{\bf 所需条件：} $\Od$在$\Hom(\pi_1, \Od)$上的共轭作用的几何不变量理论（GIT）商分析。
\end{problem}

\begin{problem}[体-边界对应的全息对偶]\label{prob:holographic}
定理\ref{thm:cercis_bulk_boundary}将Cercis解释为边界到体态的距离。
是否存在更深层的``全息对偶''（holographic duality）？
即：边界的Cercis分布是否完全决定了体态的某些全局几何量？
{\bf 所需条件：} 在2-复形上建立类似于AdS/CFT的离散全息原理。
目前这是推测性的——离散全息对偶仅在SYK模型等特定系统中被严格建立\cite{sachdev2010}。
\end{problem}

\newpage

% ══════════════════════════════════════════════════════════════════════
% SECTION 7: CONCLUSION
% ══════════════════════════════════════════════════════════════════════
\section{结论 / Conclusion}
\label{sec:conclusion}

\subsection{本文的贡献 / Contributions of This Paper}

本文在三个领域完善了SCX多专家规范理论的形式化：

\begin{enumerate}[label=(\arabic*)]
    \item \textbf{$\Od$格点规范理论（第\ref{sec:od}节）：}
    建立了完整的非阿贝尔离散规范理论——$\Od$值边联络、面Wilson圈曲率、
    非阿贝尔规范变换、$\Od$扭曲能量泛函。
    证明了小联络下规范固定退化为$\so(d)$上的标准图霍奇分解（定理\ref{thm:od_hodge_fix}），
    大扭曲下解的存在性由紧致性保证。
    这使SCX框架首次能够严格处理专家的旋转规范自由度（ACE势比较的核心需求）。

    \item \textbf{Dijkgraaf-Witten离散TQFT（第\ref{sec:dw}节）：}
    将专家图构造为2-复形，在其上定义DW离散TQFT。
    证明了SCX 2-复形同伦于$K_M$，第一Betti数为$\ZBetti = (M-1)(M-2)/2$（定理\ref{thm:scx_homotopy}）。
    建立了DW配分函数$Z_{\DW}$与平坦联络模空间$\Gflatquot$的精确关系（定理\ref{thm:dw_partition}）。
    核心成果：{\bf 证明了Cercis的调和分量$r_{\text{harm}}$是平坦联络模空间$\Gflatquot$在平凡联络处的切向量}
    （定理\ref{thm:dw_cercis_correspondence}），从而赋予调和残差以精确的模空间几何意义。

    \item \textbf{信息几何体-边界对应（第\ref{sec:info}节）：}
    在$\Situs$流形上引入Fisher信息度量，将专家输出建模为指数族。
    证明了指数族下Fisher测地线长度与KL散度的局部等价性（定理\ref{thm:fisher_kl_equivalence}）。
    建立了Cercis分数的信息几何解释：
    $\cercis = \min_{\text{纯规范体态}} (\text{边界到体态的Fisher测地距离})^2$
    $\approx 2\min_{\text{纯规范体态}} \KL(P_{\text{obs}} \| P_{\text{pure gauge}})$
    （定理\ref{thm:cercis_bulk_boundary}）。
\end{enumerate}

\subsection{数学地位 / Mathematical Status}

本文中出现的所有断言的分类：
\begin{itemize}
    \item \textbf{定义：} 15个（\ref{def:od_gauge}, \ref{def:2complex_od}, \ref{def:od_flat}, \ref{def:od_energy}, \ref{def:od_cercis}, \ref{def:scx_2complex}, \ref{def:dw_connection}, \ref{def:situs_stat}, \ref{def:fisher_metric}, \ref{def:exp_fam}, \ref{def:kl}, \ref{def:bulk_boundary}）
    \item \textbf{定理（含证明）：} 9个（\ref{thm:od_flat_char}, \ref{thm:od_hodge_fix}, \ref{thm:od_cercis_inv}, \ref{thm:scx_homotopy}, \ref{thm:dw_partition}, \ref{thm:dw_cercis_correspondence}, \ref{thm:harmonic_moduli_geometry}, \ref{thm:fisher_kl_equivalence}, \ref{thm:cercis_bulk_boundary}）
    \item \textbf{引理（含证明）：} 2个（\ref{lem:cartan}, \ref{lem:kl_bregman}）
    \item \textbf{命题（含证明）：} 4个（\ref{prop:od_stabilizer}, \ref{prop:wilson_gauge_od}, \ref{prop:homology}（未显示），\ref{prop:homology}）
    \item \textbf{推论：} 4个（\ref{cor:betti_values}, \ref{cor:dw_abelian}, \ref{cor:cercis_kl}）
    \item \textbf{未解决问题：} 5个（\ref{prob:moduli_space}--\ref{prob:holographic}）
\end{itemize}

没有任何断言是``类比''或``启发式''的。每个定理都有明确的假设和完整的证明（或证明概要）。
未解决的问题被诚实标记。

\subsection{应用前景与推广方向 / Applications and Generalization Directions}

\begin{enumerate}[label=(\arabic*)]
    \item \textbf{$\Od$规范固定在ACE势能面审计中的应用。}
    当比较多个ACE势的预测时，$\Od$ Cercis分数$\cercis_{\Od}$直接量化了
    不同ACE基函数主轴选择导致的不一致性。在材料科学的高通量筛选管线中，
    该分数可用于自动标记``旋转不一致''的专家对，触发ACE基函数的重新正交化。

    \item \textbf{DW配分函数用于专家系统的复杂度分析。}
    $\ZBetti = (M-1)(M-2)/2$提供了一个\emph{先验}的复杂度指标：
    在运行任何规范固定之前，即可预测系统潜在的``不一致模式''数量。
    对于$M$个专家的系统，审计复杂度以$O(M^2)$增长（每个专家对需比较），
    但不一致模式数以$O(M^2)$增长（$\ZBetti \sim M^2/2$）。
    这意味着：随着专家数增加，不一致模式的数量以与审计成本相同的渐近速率增长——
    系统不会因规模扩大而变得更不可审计。

    \item \textbf{信息几何Cercis用于分布型专家（如贝叶斯神经网络）。}
    对于输出完整后验分布的专家（而非点估计），
    KL正则化Cercis（定理\ref{thm:kl_regularized_cercis}）提供了自然的比较度量。
    此时``边联络''$A_e$为两个分布之间的最优传输映射，
    Cercis为所有边上最优传输代价的规范固定残差之和。

    \item \textbf{三维ACE势的$\SOthree$规范固定数值实验。}
    对于$d=3$，$\SOthree$的Cartan半径为$\pi$。
    在此半径内的ACE势比较可使用算法\ref{def:od_algorithm}直接处理。
    对于超出Cartan半径的大旋转（如ACE主轴反转），
    需初始化多个起始点并用全局优化（如模拟退火）寻找全局最小值。
    该场景对应定理\ref{thm:od_hodge_fix}中描述的``Gribov模糊''分支结构。

    \item \textbf{与持续同调工具的整合。}
    DW TQFT视角下，$\ZBetti$不仅是模空间的维数，也是$H_1$持续性条形码中
    ``长条''的个数。当SCX图的构型数$N$增大时，通过持续同调\cite{edelsbrunner2008}
    可区分``真正的调和模式''（跨所有$N$持续存在）与``有限尺寸效应''（仅跨部分构型存在）。
    Cercis的调和分量按持续性加权后可获得更强的鲁棒性。

    \item \textbf{高维输出空间（$d > 3$）的处理。}
    对于$d > 3$，$\Od$的维数为$d(d-1)/2$，Cartan半径缩小。
    在小联络假设更可能成立的条件下，算法\ref{def:od_algorithm}的线性近似更精确。
    但对于高维空间，$\Od$的几何复杂性（大圆距离、测地线的非唯一性）
    可能导致实际收敛困难，需采用$\Od$上的对数映射的稳定数值实现。

    \item \textbf{与$\Situs$物理位置编码的联合。}
    当$\Situs$的Fisher度量退化为欧氏度量（对应于正态分布族），
    信息几何Cercis（定理\ref{thm:cercis_bulk_boundary}）退化为fiber\_bundle.tex的欧氏Cercis。
    二者的差异直接度量了专家输出分布的非高斯性——
    一个可独立验证的统计量。
\end{enumerate}

\subsection{具体计算示例：三专家ACE势能面比较 / Worked Example: Three-Expert ACE PES Comparison}

为阐明本文的形式化框架的实际应用，考虑以下具体场景：

{\bf 设定：} $M=3$个专家（ACE势），$N=100$个构型（分子动力学轨迹的快照），
输出空间$d=3$（能量、力x分量、力y分量）。每个专家返回一个3维向量。

{\bf 步骤1 — 构建SCX 2-复形：}
顶点数$|\verts| = 300$，边数$|\edgs| = 3(99) + 100 \times 3 \times 2 = 297 + 600 = 897$，
面数$|\faces| = 99 \times \binom{3}{2} = 99 \times 3 = 297$。
第一Betti数$\ZBetti = (3-1)(3-2)/2 = 1$——恰好一个独立的不一致模式。

{\bf 步骤2 — 边赋值：}
参数边$a_{(k,m) \to (k+1,m)} = x_m^{k+1} - x_m^k \in \R^3$（平移分量），
或旋转分量$R_{(k,m) \to (k+1,m)} \in \SOthree$（若考虑旋转规范）。
专家边$a_{(k,i) \to (k,j)} = x_i^k - x_j^k$（或对应的旋转）。

{\bf 步骤3 — 规范固定：}
平移规范：构建$B \in \R^{897 \times 300}$，求解$B^T B g = B^T A$（3个坐标分别求解），
带零模固定$\sum_v g_v = 0$。
得到Cercis平移分数$\cercis_{\R^3}$。

旋转规范（如果适用）：运行算法\ref{def:od_algorithm}在$\SOthree^{300}$上。
得到Cercis旋转分数$\cercis_{\SOthree}$。
总Cercis：$\cercis_{\SEN} = \cercis_{\R^3} + \cercis_{\SOthree}$。

{\bf 步骤4 — 调和分析：}
Cercis调和分量$r_{\text{harm}}$的维数：
\begin{itemize}
    \item 平移部分：$3 \times 1 = 3$维（3个坐标方向，1个独立回路模式）；
    \item 旋转部分：$3 \times 1 = 3$维（$\mathfrak{so}(3)$的3个基，1个独立回路模式）。
\end{itemize}
总计6个独立的调和方向——6种本质上不同的``专家不一致方式''。

{\bf 步骤5 — DW解释：}
若将位移离散化（如四舍五入到$\Z_{10}$格点），DW配分函数为
$Z_{\DW}(\mathcal{K}_{\text{SCX}}, \Z_{10}) = 10^1 = 10$（平移）或
$Z_{\DW} = |\SOthree_{\text{离散}}|^1$（旋转，取决于离散化精度）。
这10个（或更多）状态对应10种规范不等价的全局不一致模式。

{\bf 步骤6 — 信息几何解释：}
将每个构型$k$处每个专家$m$的输出建模为$\mathcal{N}(x_m^k, \Sigma_m^k)$（正态分布）。
Fisher度量退化为Mahalanobis距离。Cercis等值于从观测分布到``完美对齐''分布
的最小Mahalanobis距离平方和。
\end{enumerate}

\subsection{最终评注 / Final Remark}

SCX多专家系统的规范结构最终归结为三个相互关联的数学对象：
\begin{enumerate}[label=(\roman*)]
    \item 专家比较图上的非阿贝尔$\Od$格点规范场（描述旋转自由度）；
    \item 2-复形上的Dijkgraaf-Witten TQFT（描述平坦联络模空间，解释调和分量）；
    \item $\Situs$信息流形上的Fisher-Rao度量（提供Cercis的概率解释和体-边界对应）。
\end{enumerate}

这三个结构分别回答了以下问题：
\begin{itemize}
    \item ``专家的旋转参考标架如何被全局对齐？'' — $\Od$格点规范固定
    \item ``专家系统有多少种本质上不同的不一致方式？'' — DW模空间$\Gflatquot$，维数$(M-1)(M-2)/2$
    \item ``Cercis分数在概率上意味着什么？'' — 到全局一致分布的最小信息距离
\end{itemize}

我们相信，这一形式化为SCX的规范理论提供了坚实的数学基础，
每一个断言都有明确的定义和可验证的数学地位。

\newpage

% ══════════════════════════════════════════════════════════════════════
% BIBLIOGRAPHY
% ══════════════════════════════════════════════════════════════════════
\begin{thebibliography}{99}

\bibitem{scx_fiber_bundle}
SCX Theory Architecture Group.
``势能面不齐的离散几何：SCX规范理论的图霍奇形式化 / Discrete Geometry of PES Misalignment:
A Graph Hodge Formalization of SCX Gauge Theory.''
Working paper, 2026.

\bibitem{dijkgraaf1990}
R.~Dijkgraaf and E.~Witten.
``Topological Gauge Theories and Group Cohomology.''
\emph{Communications in Mathematical Physics}, 129(2):393--429, 1990.

\bibitem{witten1991}
E.~Witten.
``On quantum gauge theories in two dimensions.''
\emph{Communications in Mathematical Physics}, 141(1):153--209, 1991.

\bibitem{freed1995}
D.~S.~Freed and F.~Quinn.
``Chern-Simons theory with finite gauge group.''
\emph{Communications in Mathematical Physics}, 156(3):435--472, 1993.

\bibitem{amari2016}
S.~Amari.
\emph{Information Geometry and Its Applications}.
Springer, 2016.

\bibitem{ay2017}
N.~Ay, J.~Jost, H.~V\^an L\^e, and L.~Schwachh\"ofer.
\emph{Information Geometry}.
Springer, 2017.

\bibitem{rao1945}
C.~R.~Rao.
``Information and accuracy attainable in the estimation of statistical parameters.''
\emph{Bulletin of the Calcutta Mathematical Society}, 37:81--91, 1945.

\bibitem{cover2006}
T.~M.~Cover and J.~A.~Thomas.
\emph{Elements of Information Theory}, 2nd ed.
Wiley, 2006.

\bibitem{lim2020}
L.-H.~Lim.
``Hodge Laplacians on Graphs.''
\emph{SIAM Review}, 62(3):685--715, 2020.

\bibitem{jiang2011}
X.~Jiang, L.-H.~Lim, Y.~Yao, and Y.~Ye.
``Statistical ranking and combinatorial Hodge theory.''
\emph{Mathematical Programming}, 127(1):203--244, 2011.

\bibitem{hatcher2002}
A.~Hatcher.
\emph{Algebraic Topology}.
Cambridge University Press, 2002.

\bibitem{bott1982}
R.~Bott and L.~Tu.
\emph{Differential Forms in Algebraic Topology}.
Springer, 1982.

\bibitem{kobayashi1963}
S.~Kobayashi and K.~Nomizu.
\emph{Foundations of Differential Geometry, Vol.~I}.
Wiley-Interscience, 1963.

\bibitem{nakahara2003}
M.~Nakahara.
\emph{Geometry, Topology and Physics}, 2nd ed.
Taylor \& Francis, 2003.

\bibitem{chung1997}
F.~R.~K.~Chung.
\emph{Spectral Graph Theory}.
American Mathematical Society, 1997.

\bibitem{desbrun2005}
M.~Desbrun, A.~N.~Hirani, M.~Leok, and J.~E.~Marsden.
``Discrete Exterior Calculus.''
\emph{arXiv:math/0508341}, 2005.

\bibitem{grady2011}
L.~J.~Grady and J.~R.~Polimeni.
\emph{Discrete Calculus: Applied Analysis on Graphs for Computational Science}.
Springer, 2011.

\bibitem{mednykh}
A.~D.~Mednykh.
``Determination of the number of nonequivalent coverings over a compact Riemann surface.''
\emph{Soviet Mathematics Doklady}, 19:318--320, 1978.

\bibitem{sachdev2010}
S.~Sachdev.
``Holographic metals and the fractionalized Fermi liquid.''
\emph{Physical Review Letters}, 105(15):151602, 2010.

\bibitem{montvay1994}
I.~Montvay and G.~M\"unster.
\emph{Quantum Fields on a Lattice}.
Cambridge University Press, 1994.

\bibitem{creutz1983}
M.~Creutz.
\emph{Quarks, Gluons and Lattices}.
Cambridge University Press, 1983.

\bibitem{wilson1974}
K.~G.~Wilson.
``Confinement of quarks.''
\emph{Physical Review D}, 10(8):2445--2459, 1974.

\bibitem{schaub2020}
M.~T.~Schaub, A.~R.~Benson, P.~Horn, G.~Lippner, and A.~Jadbabaie.
``Random Walks on Simplicial Complexes and the normalized Hodge 1-Laplacian.''
\emph{SIAM Review}, 62(2):353--391, 2020.

\bibitem{nocedal2006}
J.~Nocedal and S.~J.~Wright.
\emph{Numerical Optimization}, 2nd ed.
Springer, 2006.

\bibitem{tishby2000}
N.~Tishby, F.~C.~Pereira, and W.~Bialek.
``The information bottleneck method.''
\emph{Proceedings of the 37th Annual Allerton Conference}, 368--377, 1999.

\bibitem{edelsbrunner2008}
H.~Edelsbrunner and J.~Harer.
\emph{Computational Topology: An Introduction}.
American Mathematical Society, 2010.

\end{thebibliography}

\end{document}